\documentclass[11pt,a4paper]{article}

% ========= ПАКЕТЫ (без inputenc/fontenc – они не нужны для LuaLaTeX) =========
\usepackage[english, bidi=default, layout=counters]{babel}
\usepackage{amsmath,amssymb,amsfonts}
\usepackage{geometry}
\geometry{left=1.25cm,right=1.25cm,top=1.25cm,bottom=3.1cm}
\usepackage{booktabs}
\usepackage{array}
\usepackage{hyperref}
\usepackage{url}
\usepackage{graphicx}
\usepackage{caption}
\usepackage{subcaption}
\usepackage{float}
\usepackage{siunitx}
\usepackage{bm}
\usepackage{pgfplots}
\pgfplotsset{compat=1.18}
\usepackage{xcolor}
\usepackage{titlesec}
\usepackage{fancyhdr}
\usepackage{microtype}
\usepackage{fontspec}
\usepackage{tikz}
\usetikzlibrary{shapes,arrows,arrows.meta,positioning,decorations.pathmorphing,patterns,calc}
\usepackage{multicol}
\usepackage{textcomp}

\tikzset{>=stealth}

% ========= ЯЗЫКИ =========
\babelprovide[import, main, maparabic, alph=alph]{arabic}
\babelprovide[import, onchar=ids fonts]{english}

% ========= АРАБСКИЕ ШРИФТЫ (Amiri — полная поддержка) =========
\defaultfontfeatures{Ligatures=TeX}
\setmainfont{Amiri}[
Script=Arabic,
Mapping=arabicdigits,
Ligatures=TeX,
BoldFont=Amiri-Bold,
ItalicFont=Amiri-Italic,
BoldItalicFont=Amiri-BoldItalic
]
\setsansfont{Amiri}[
Script=Arabic,
Mapping=arabicdigits
]
\newfontfamily{\arabicfonttt}{Amiri}[
Script=Arabic,
Mapping=arabicdigits
]

% Шрифты для латинского текста (английского)
\babelfont[english]{rm}{Times New Roman}
\babelfont[english]{sf}{Arial}
\babelfont[english]{tt}{Courier New}

% ========= ГЛОБАЛЬНАЯ НАСТРОЙКА ДЛЯ ВСЕХ РИСУНКОВ =========
% Принудительно переключаем направление внутри tikzpicture на LTR,
% чтобы координаты не зеркалировались, но арабский текст оставался RTL.
\tikzset{every picture/.style={execute at begin picture={\textdir TLT}}}

% ========= ЦВЕТА =========
\definecolor{skyblue}{RGB}{0,102,204}
\definecolor{darkblue}{RGB}{0,0,139}
\definecolor{gold}{RGB}{255,215,0}

% ========= КОМАНДЫ =========
\newcommand{\eps}{\varepsilon}
\newcommand{\kappac}{\kappa_c}
\newcommand{\Keff}{K_{\mathrm{eff}}}
\newcommand{\betaf}{\beta}
\newcommand{\df}{d_f}

% ========= ЗАГОЛОВКИ =========
\titleformat{\section}{\LARGE\bfseries\color{darkblue}}{\thesection}{1em}{}
\titleformat{\subsection}{\normalsize\bfseries\color{darkblue}}{\thesubsection}{1em}{}
\titleformat{\subsubsection}{\normalsize\itshape}{\thesubsubsection}{1em}{}

% ========= КОЛОНТИТУЛЫ =========
\fancypagestyle{firstpage}{%
	\fancyhf{}%
	\renewcommand{\headrulewidth}{0pt}%
	\renewcommand{\footrulewidth}{0pt}%
	\fancyfoot[C]{%
		\begin{minipage}[t]{\textwidth}
			\centering
			\textcolor{gold}{\rule{\textwidth}{1.6pt}}\\[5pt]
			{\small \textsuperscript{1}أبحاث ياكوشيف، \url{https://yuct.org/}} \hfill
			{\small بروتوكول YPSDC: \url{https://ypsdc.com/}}\\[-2pt]
			\textcolor{gold}{\rule{\textwidth}{1.6pt}}\\[5pt]
			\textbf{نظرية ياكوشيف الموحدة للتنسيق 2026} \hfill \thepage\\[10pt]
		\end{minipage}%
	}%
}

\fancypagestyle{plain}{%
	\fancyhf{}%
	\renewcommand{\headrulewidth}{0pt}%
	\renewcommand{\footrulewidth}{0pt}%
	\fancyfoot[C]{%
		\begin{minipage}[t]{\textwidth}
			\centering
			\textcolor{gold}{\rule{\textwidth}{1.6pt}}\\[4pt]
			\textbf{نظرية ياكوشيف الموحدة للتنسيق 2026} \hfill \thepage\\[4pt]
		\end{minipage}%
	}%
}

\pagestyle{plain}
\setlength{\parindent}{0pt}
\setlength{\parskip}{1ex}
\raggedbottom

\hypersetup{
	colorlinks=true,
	linkcolor=blue,
	citecolor=blue,
	urlcolor=blue,
}

% ========= ЗАГОЛОВОК СТАТЬИ =========
\newcommand{\makeheader}{%
	\noindent
	\begin{minipage}{0.17\textwidth}
		\raggedright
		{\sffamily\bfseries\color{skyblue}\fontsize{40}{40}\selectfont YUCT}%
	\end{minipage}%
	\hspace{30pt}
	\begin{minipage}{0.32\textwidth}
		\raggedright
		{\sffamily\fontsize{11}{13}\selectfont نظرية ياكوشيف\\ الموحدة\\ للتنسيق}%
	\end{minipage}%
	\hfill
	\begin{minipage}{0.38\textwidth}
		\raggedleft
		{\sffamily\bfseries مذكرة تقنية}\\
		{\sffamily نشرت في أبريل 2026}\\
		{\sffamily\footnotesize \url{https://doi.org/10.5281/zenodo.18444598}}%
	\end{minipage}
	\par\vspace{5pt}
	\noindent\textcolor{gold}{\rule{\textwidth}{1.6pt}}
	\par\vspace{0pt}
}

\begin{document}
	\thispagestyle{firstpage}
	\makeheader
	
	\vspace{0cm}
	{\LARGE\bfseries الكارثية الجاذبية: نموذج للاصطدامات الدورية لجرم عابر هائل بنظام الأرض–القمر والمحيط الحيوي \par}
	\vspace{0.2cm}
	
	{\large أليكسي ف. ياكوشيف\textsuperscript{1}} \\
	\vspace{0.0cm}
	
	{\color{darkblue}\textbf{\LARGE المستخلص}} \par
	{\large
		\noindent
		تقدم هذه الورقة نموذجاً موحداً للانقراضات الجماعية الدورية والكوارث الأرضية ونشأة الحياة في إطار نظرية التنسيق الموحدة لياكوشيف (YUCT). نبين أن جرمًا عابرًا هائلًا - بقايا قمر ثان للأرض - بعد تفاعل ثلاثي الأجسام فوضوي مع نظام الأرض–القمر، قُذف في مدار بيضاوي شديد الانحراف بزمن دوري قدره \(26.1\) مليون سنة. تمريراته المتكررة عبر سحابة أورت الداخلية تحفز زخات المذنبات، بينما عبور نفس مستوى المجرة الذي يزعزع استقرار سحابة أورت يقلل أيضاً من كفاءة التنسيق \(K_{\mathrm{eff}}\) لدثار الأرض. هذا الانخفاض يخفض لزوجة الدثار والاحتكاك على طول مناطق الصدوع، منتجاً شلالاً متزامناً: طفرات في الاصطدامات، براكين الفيضانات البازلتية، تجوالات مغناطيسية أرضية، تسارع التصدع/الاصطدام، انقطاع الأكسجين المحيطي، والانقراضات الجماعية - كلها تشترك في نفس الدورة \(26\)–\(30\) مليون سنة. نشتق الزمن الدوري مباشرة من الأس YUCT العام \(\beta = 2/3\) ونحدد كمياً كل خطوة في الشلال بصيغ قابلة للاختبار. يشرح النموذج الازدواجية القمرية، كتلة مياه الأرض، والنبض الجيولوجي البالغ 27.5 مليون سنة الموثق بشكل مستقل بواسطة رامبينو وآخرين. نجادل أيضاً بأن الاصطدام عقّم نظام الأرض–القمر (عبر غلاف جوي من بخار ساخن حامضي للغاية)، رافضاً نظرية البانسبيرميا، مما يعني ضمناً نشأة حيوية محلية في الهاديان. نقدم ثلاث تنبؤات قابلة للتكذيب: (1) سينتج العبور القادم لمستوى المجرة تجوالاً مغناطيسياً أرضياً يدوم \(440 \pm 80\) سنة، (2) سيزيد معدل اصطدام الأجسام دون الكيلومترية بعامل \(10^{2}\)–\(10^{3}\) لمدة 1–2 مليون سنة، (3) ستبدأ براكين الفيضانات البازلتية وانقطاع الأكسجين المحيطي في غضون \(\pm 0.5\) مليون سنة من طفرة الاصطدام. وهكذا يوحد شلال YUCT الفيزياء الفلكية، الجيوفيزياء، علم الأحياء القديمة، وعلم الأحياء الفلكي تحت إطار كمي واحد.}
	
	\vspace{0.1cm}
	\noindent
	{\color{darkblue}\textbf{الكلمات المفتاحية:}} الكارثية الجاذبية، الانقراضات الجماعية، الفوهات الصدمية، القمر الثاني للأرض، الماء على الأرض، YUCT، كفاءة التنسيق، المد والجزر المجري، دورية 26 مليون سنة.
	
	\vspace{0.1cm}
	\noindent
	{\color{darkblue}\textbf{الاقتباس:}} Yakushev AV. Gravitational Catastrophism: A Model of Periodic Impacts of a Massive Transient Object on the Earth–Moon System and the Biosphere. 2026. DOI: \url{https://doi.org/10.5281/zenodo.18444598} (this volume).
	
	\vspace{0.4cm}
	
	% FIGURE 1: Земля и две луны (обёрнута в english для правильного центрирования)
	\begin{otherlanguage}{english}
		\begin{figure*}[htbp]
			\centering
			\begin{tikzpicture}[scale=0.9, every node/.style={font=\small}]
				\shade[ball color=blue!40!white, opacity=0.8] (0,0) circle (1.2);
				\node at (0,-1.5) {الأرض};
				\shade[ball color=gray!60!white] (3.5,0) circle (0.4);
				\node at (3.5,-0.8) {القمر};
				\shade[ball color=brown!70!black] (2.2,2.2) circle (0.25);
				\node at (2.2,2.8) {القمر الثاني};
				\draw[dashed, blue!70] (0,0) ellipse (4.5 and 4.5);
				\node at (4.5,0.5) {$L_4$};
				\node at (-4.5,-0.5) {$L_5$};
				\draw[->, thick, red, decorate, decoration={snake,amplitude=0.5mm,segment length=2mm}] (2.2,2.2) -- (3.5,0.4);
				\node[red] at (3.2,1.8) {اصطدام بطيء};
				\draw[thin, white!80!black] (0,0) circle (3.5);
			\end{tikzpicture}
			\caption{إعادة بناء نظام الأرض–القمر–قمر ثان القديم. كان القمر الثاني (كتلته \(\approx 1/27\) من القمر) يقع بالقرب من نقطة لاغرانج \(L_4\) أو \(L_5\) ثم اندمج مع القمر في اصطدام منخفض السرعة، مفسراً الازدواجية القمرية ومقدماً الماء إلى الأرض.}
			\label{fig:two-moons}
		\end{figure*}
	\end{otherlanguage}
	
	\begin{multicols}{2}
		
		\section{المقدمة}
		لطالما افتقرت ظاهرة الكوارث العالمية المتكررة المسجلة في السجل الجيولوجي (الانقراضات الجماعية، ذروات أحداث الاصطدام، النشاط التكتوني) إلى تفسير موحد. يطور العمل الحالي فرضية تربط هذه الأحداث بالاقتراب الدوري لجرم عابر هائل. كان هذا الجسم موجوداً في البداية كقمر ثان للأرض؛ بعد اصطدامه بالقمر، انتقل إلى مدار ممدد للغاية يعبر بانتظام النظام الشمسي الداخلي. يؤدي الاضطراب الجاذبي الناتج عن هذا الجسم إلى إطلاق "زخات" من المذنبات والكويكبات موجهة نحو الأرض، منتجاً بذلك الدورية المرصودة للكوارث.
		
		إضافة مهمة هي التفسير في إطار نظرية التنسيق الموحدة لياكوشيف (YUCT)، التي تعتبر حقل الجاذبية المجري "قاموساً" خارجياً، والتغيرات الدورية في كفاءة التنسيق \(K_{\mathrm{eff}}\) كآلية عالمية تربط المقاييس الصغرى والكبرى.
		
		\section{الأساس الرصدي للدورية}
		\subsection{الانقراضات الجماعية وأحداث الاصطدام}
		يُظهر تحليل السجل الأحفوري على مدى الـ 250 مليون سنة الماضية \cite{Raup1984, Melott2010} دورة ذات دلالة إحصائية للانقراضات الجماعية بزمن دوري:
		\[
		T_{\text{ext}} = 26.0 \pm 0.5\ \text{مليون سنة}.
		\]
		تعطي البيانات المنقحة \(T_{\text{ext}} = 30.0 \pm 1.0\) مليون سنة. توجد دورية مماثلة لفوهات الصدم الكبيرة (قطر \(>35\) كم) — \(T_{\text{cr}} = 30.0 \pm 5.0\) مليون سنة \cite{RampinoStothers1984}، ولـ"النَّبْض الجيولوجي" — ذروات النشاط البركاني والتكتوني (\(T_{\text{geo}} = 33.0 \pm 3.0\) مليون سنة). تشير سلاسل البيانات المستقلة الثلاث هذه إلى آلية خارجية واحدة.
		
		\subsection{أكبر فوهات الصدم على الأرض}
		\begin{table}[htbp]
			\centering
			\caption{أكبر فوهات الصدم المؤكدة على الأرض.}
			\label{tab:craters}
			\footnotesize
			\begin{tabular}{lccr}
				\toprule
				\textbf{الفوهة} & \textbf{الموقع} & \textbf{القطر (كم)} & \textbf{العمر (مليون سنة)} \\
				\midrule
				فريديفورت & جنوب أفريقيا & 250–300 & 2023 \\
				سدبري & كندا & ~250 & 1850 \\
				شيكشولوب & المكسيك & ~180 & 66.5 \\
				بوبيغاي & روسيا & ~100 & 35.7 \\
				مانيكواجان & كندا & ~100 & 214 \\
				\bottomrule
			\end{tabular}
		\end{table}
		
		ذات أهمية خاصة فوهة شيكشولوب (66.5 مليون سنة)، المرتبطة بانقراض العصر الطباشيري–الباليوجيني، وفوهة بوبيغاي (35.7 مليون سنة)، التي تقع أيضاً ضمن دورة الـ26 مليون سنة.
		
		\subsection{تناقص الشدة بمرور الوقت}
		تناقصت الشدة الخلفية للانقراضات خطياً طوال الدهر الظاهري. كانت الطاقة المنطلقة من أحداث الصدم أعلى بشكل منهجي في العصور القديمة. بينما تبقى الدورية، يتناقص سعة الدورة، وهو ما يفسر على أنه تَبَاعُد تدريجي لمصدر الاضطراب عن الأرض (زيادة المحور شبه الرئيسي لمدار الجسم).
		
		\section{فرضية القمر الثاني والاصطدام بالقمر}
		\subsection{الازدواجية القمرية}
		تظهر البيانات الحديثة عن تضاريس القمر اختلافاً حاداً بين الجانب القريب والبعيد: يهيمن على الجانب القريب قشرة رقيقة وسهول حمم بركانية (ماريا)، بينما يتميز الجانب البعيد بقشرة سميكة وتضاريس جبلية. تفسر هذه الازدواجية بنموذج اصطدام الأرض بقمرين \cite{JutziAsphaug2011}. وفقاً للنموذج، كان قمر ثان (كتلته \(\sim 1/27\) من القمر الحالي) موجوداً في إحدى نقطتي لاغرانج \(L_4\) أو \(L_5\) واصطدم ببطء بالقمر، منتشراً على سطحه.
		
		\subsection{حساب كتلة القمر الثاني}
		بافتراض أن قطر القمر الثاني يساوي \(1/3\) قطر القمر وبنفس الكثافة، نحصل على:
		{\small
			\[
			M_{\text{moonlet}} = \left(\frac{1}{3}\right)^3 M_{\text{moon}} = \frac{1}{27} \cdot 7.35 \times 10^{22}\ \text{kg} \approx 2.72 \times 10^{21}\ \text{kg}.
			\]}
		
		\subsection{الزخم الزاوي للنظام}
		الزخم الزاوي الكلي لنظام الأرض–القمر هو:
		\[
		L_{\text{total}} = L_{\text{orbit}} + L_{\text{spin}} \approx 3.61 \times 10^{34}\ \text{kg}\cdot\text{m}^2/\text{s},
		\]
		80\% منه في الحركة المدارية للقمر. لا بد أن اصطدام القمر الثاني بالقمر حدث بسرعة:
		\[
		v_{\text{impact}} < 2.5\ \text{km/s},
		\]
		وإلا لحدث تدمير بدلاً من الاندماج. هذه السرعة المنخفضة ممكنة فقط عندما تتحرك الأجسام في نفس المستوى المداري بسرعات متشابهة.
		
		% =========================================================================
		% GRAVITATIONAL PRESSURE AND CHAOTIC THREE‑BODY DYNAMICS
		% =========================================================================
		\subsection{الضغط الجاذبي والديناميكيات الفوضوية للأجسام الثلاثة}
		إن وجود قمرين هائلين (القمر والقمر الثاني) حول الأرض البدائية خلق حقل مد وجزر قوي ومتغير زمنياً بشكل استثنائي. من المحتمل أن الضغط الجاذبي الذي مارسه هذان الجسمان على الغلاف الصخري للأرض، إلى جانب الفوضى الجوهرية لنظام الأجسام الثلاثة، لعب دوراً حاسماً في تكسير القشرة البدائية وإطلاق الصفائح التكتونية.
		
		\subsubsection{تقدير الضغط الجاذبي}
		بالنسبة لكتلة نقطية \(M\) على مسافة \(d\)، فإن تسارع المد عبر جسم نصف قطره \(R\) هو \(\Delta a = 2GM R / d^{3}\). ينتج تدرج التسارع هذا إجهاداً (ضغطاً) على الغلاف الصخري. يمكن تقدير ضغط المد المميز كما يلي:
		\[
		P_{\text{tidal}} \sim \frac{GM \rho R}{d^{2}} \cdot \frac{R}{d} = \frac{GM \rho R^{2}}{d^{3}},
		\]
		حيث \(\rho \approx 3300\;\text{kg/m}^{3}\) هو متوسط كثافة القشرة الأرضية.
		
		بأخذ مسافة القمر المبكرة \(d_{\text{Moon}} \approx 3\times10^{8}\;\text{m}\) (حوالي نصف المسافة الحالية) والقمر الثاني على مسافة مماثلة أو أكبر قليلاً (مثل \(d_{2} \approx 4\times10^{8}\;\text{m}\)) بكتلة \(M_{2} \approx 2.7\times10^{21}\;\text{kg}\)، نحصل على:
		\[
		\begin{aligned}
			P_{\text{Moon}} &\sim \frac{6.67\times10^{-11} \cdot 7.35\times10^{22} \cdot 3300 \cdot (6.37\times10^{6})^{2}}{(3\times10^{8})^{3}} \\
			&\approx 3.5 \times 10^{4}\;\text{Pa} \;(0.035\;\text{MPa}),\\
			P_{\text{moonlet}} &\sim \frac{6.67\times10^{-11} \cdot 2.7\times10^{21} \cdot 3300 \cdot (6.37\times10^{6})^{2}}{(4\times10^{8})^{3}} \\
			&\approx 4.7 \times 10^{3}\;\text{Pa} \;(0.0047\;\text{MPa}).
		\end{aligned}
		\]
		
		على الرغم من أن هذه القيم صغيرة مقارنة بمقاومة الصخور السليمة (عشرات MPa)، فإن النقطة الحاسمة هي أنها **متغيرة زمنياً** و**غير ثابتة**. تحرك القمران في مدارات غير دائرية، وتسبب تفاعلهما الجاذبي المتبادل في تغيرات مستمرة في اتجاه وحجم قوة المد. نتج عن ذلك سعة إجهاد دورية تصل إلى عدة عشرات من الكيلوباسكالات – كافية لبدء تشقق التعب على مدى أزمنة جيولوجية، خاصة بوجود نقاط ضعف سابقة.
		
		\subsubsection{الديناميكيات الفوضوية للأجسام الثلاثة}
		نظام الأرض–القمر–قمر صغير هو مثال كلاسيكي لمسألة الأجسام الثلاثة الهرمية. تُظهر التكاملات العددية لمثل هذه الأنظمة (مثل مسألة الأجسام الثلاثة المقيدة) أن المدارات في جوار نقطتي لاغرانج \(L_{4}\) و \(L_{5}\) ليست مستقرة تماماً عندما يكون للجسم الثالث كتلة غير مهملة. تؤدي الاضطرابات الصغيرة (من مد وجزر الشمس، اللامركزية، أو الكواكب الأخرى) إلى حركة فوضوية: يمكن أن يخضع مدار القمر الصغير لتقلبات كبيرة لا يمكن التنبؤ بها في اللامركزية والميل قبل اصطدام نهائي بالقمر أو قذفه.
		
		لهذا النظام الفوضوي نتيجتان مهمتان:
		\begin{enumerate}
			\item \textbf{قوة دفع مدية شديدة التغير}. كلما اقترب القمر الصغير من الأرض في مدار لامركزي، يمكن أن يزداد الضغط المدى مؤقتاً بمقدار مرتبة أو اثنتين. يمكن للحلقات القصيرة من الإجهاد العالي (تصل إلى عدة MPa) أن تتجاوز قوة الخضوع للغلاف الصخري المبكر المكسور.
			\item \textbf{ضخ رنيني لأنماط الغلاف الصخري}. يثير الطيف الترددي العريض للمحرك الفوضوي بفعالية الأنماط الذاتية للأرض الصلبة. يمكن للتعزيز الرنيني لأنماط الانثناء القشري تركيز الإجهاد في مواقع محددة، مما يعزز تشكل الأودية المتصدعة.
		\end{enumerate}
		
		\subsubsection{الآثار المترتبة على تمزق الغلاف الصخري وبدء الصفائح التكتونية}
		قدم الجمع بين الضغط المدى الدوري المستدام (من القمر) وأحداث الإجهاد العالي المتقطعة الفوضوية (من القمر الصغير) آلية طبيعية لـ:
		\begin{itemize}
			\item \textbf{تشقق التعب} للقشرة البدائية، مما يخلق شبكة من الشقوق العميقة.
			\item \textbf{توطين الانفعال} على طول مناطق الضعف الموجودة مسبقاً، مما يؤدي في النهاية إلى التصدع القاري.
			\item \textbf{تحفيز الاندساس} عندما ينعكس حقل الإجهاد بشكل متقطع، دافعاً كتلة قشرية تحت أخرى.
		\end{itemize}
		
		بمجرد إنشاء حدود الصفائح، فإن استمرار قوة المد من القمر (بعد اندماج القمر الصغير) يمكن أن يحافظ على حركة الصفائح، كما تشير الدراسات الحديثة التي تربط دوران الأرض ومد وجزر القمر بالحمل الحراري الوشاحي \cite{Rampino2021}. وبالتالي، ربما كانت مرحلة الأجسام الثلاثة الفوضوية لنظام الأرض–القمر هي "المفتاح" الحاسم الذي حول كوكب الغطاء الراكد إلى كوكب صفائح تكتونية.
		
		\subsubsection{نتيجة قابلة للاختبار}
		يتنبأ النموذج بأن بداية الصفائح التكتونية يجب أن ترتبط زمنياً بالمرحلة الفوضوية للتفاعل ثلاثي الأجسام، أي بين 4.4 و 4.0 مليار سنة. يمكن استخدام التأريخ عالي الدقة لأقدم الأفيوليت (مثل حزام إيسوا السبراقطي البالغ 4.0 مليار سنة) والأدلة الجيوكيميائية على الاندساس المبكر لاختبار هذا التوقع. علاوة على ذلك، فإن الدورية المميزة لأحداث التصدع التي لوحظت لاحقاً في تاريخ الأرض (نبض 27.5 مليون سنة) قد تكون من بقايا طيف الدفع الفوضوي الأصلي.
		
		\vspace{0.2cm}
		\noindent
		\textbf{ملاحظة للعمل المستقبلي:} هناك حاجة إلى محاكاة عددية كاملة للديناميكيات الفوضوية للأجسام الثلاثة مع التبدد المدى الواقعي لتحديد كمية احتمال تمزق الغلاف الصخري. هذه المحاكيات ممكنة باستخدام مقسمات الأجسام N الحديثة (مثل REBOUND، ASSIST) وستكون موضوع ورقة لاحقة.
		% =========================================================================
		
		% =========================================================================
		% THREE‑BODY CHAOS AS A CATALYST FOR COLLISION
		% =========================================================================
		\subsection{فوضى الأجسام الثلاثة وحتمية الاصطدام القمري}
		يشكل نظام الأرض–القمر–قمر صغير المبكر مشكلة الأجسام الثلاثة الهرمية. حتى إذا كان القمر الصغير يقع في البداية بالقرب من نقطتي لاغرانج المستقرتين \(L_4\) أو \(L_5\) (انظر الشكل 1)، فإن مداره ليس آمناً بشكل دائم. مسألة الأجسام الثلاثة المقيدة غير قابلة للتكامل، والاضطرابات الصغيرة من الشمس والمشتري واللامركزية غير الصفرية للقمر تدفع الانتشار الفوضوي.
		
		\subsubsection{تقديرات تحليلية لعدم الاستقرار}
		من نظرية مسألة الأجسام الثلاثة المقيدة الإهليلجية، فإن **زمن ليابونوف** للأجسام في المدارات الشرغوفية حول \(L_4\)/\(L_5\) هو من رتبة بضع مئات من الفترات المدارية للأولية (القمر). مع الفترة المدارية المبكرة للقمر \(\sim 10\) أيام، فإن زمن ليابونوف هو \(\sim 10^{3}\) يوم \(\approx 3\) سنوات – قصير للغاية. وبالتالي، حتى الاضطراب الصغير يؤدي إلى تجوال فوضوي على مقاييس زمنية أقصر بكثير من عمر النظام الشمسي.
		
		والنتيجة هي أن القمر الصغير يغادر حتمًا جوار \(L_4\)/\(L_5\) ويدخل نظاماً تصبح فيه الاقترابات القريبة من القمر ممكنة. تُظهر الدراسات العددية لأنظمة الأجسام الثلاثة المماثلة (مثل قمري زحل جانوس-إبيمثيوس \cite{JanusEpimetheus}) أن الترتيبات المدارية المشتركة عابرة على مقاييس زمنية \(10^{5}\)–\(10^{7}\) سنة. بقياس ذلك، من المتوقع أن يصبح نظام الأرض–القمر–قمر صغير غير مستقر على مقياس زمني أقصر بكثير من عمر النظام الشمسي، مما يجعل الاصطدام منخفض السرعة محتملاً إحصائياً. ستكون محاكاة مونت كارلو كاملة (مخطط لها في دراسة لاحقة) مطلوبة لتقدير الاحتمال الدقيق.
		
		\subsubsection{رسم تخطيطي للهروب الفوضوي}
		\begin{figure}[H]
			\centering
			\resizebox{\columnwidth}{!}{%
				\begin{tikzpicture}[scale=1.2, every node/.style={font=\small}]
					\shade[ball color=blue!40!white] (0,0) circle (0.6);
					\node at (0,-0.9) {الأرض};
					\draw[dashed, gray!50] (0,0) circle (2.5);
					\shade[ball color=gray!60!white] (2.5,0) circle (0.25);
					\node at (2.5,-0.5) {القمر};
					\fill[green!70!black] (1.25,2.165) circle (0.08);
					\node[green!70!black] at (1.5,2.5) {$L_4$};
					\fill[green!70!black] (1.25,-2.165) circle (0.08);
					\node[green!70!black] at (1.5,-2.5) {$L_5$};
					\shade[ball color=brown!70!black] (1.25,2.165) circle (0.12);
					\draw[red, thick, decorate, decoration={snake,amplitude=0.8mm,segment length=3mm,post length=0mm}] 
					(1.25,2.165) .. controls (1.0,1.0) and (1.5,0.5) .. (2.2,0.2);
					\draw[red, thick, dashed] (2.2,0.2) .. controls (2.5,0.0) and (2.6,-0.3) .. (2.5,-0.2);
					\node[red, right] at (2.6,0.0) {هروب فوضوي};
					\draw[->, thick, orange] (-3.5,-2) -- (-3.5,-3) node[midway, left] {اضطرابات شمسية};
				\end{tikzpicture}%
			}
			\caption{رسم تخطيطي لنظام الأجسام الثلاثة الفوضوي. يقع القمر الصغير في البداية بالقرب من نقطة لاغرانج \(L_4\) (أو \(L_5\)). تؤدي الاضطرابات الخارجية من الشمس ولامركزية القمر إلى انتشار فوضوي (سهم أحمر مموج)، مما يؤدي في النهاية إلى اقتراب قريب واصطدام بالقمر. زمن ليابونوف لمثل هذا النظام هو بضع سنوات فقط، مما يجعل الاصطدام حتمياً إحصائياً على المقاييس الزمنية الجيولوجية.}
			\label{fig:threebody}
		\end{figure}
		
		\subsubsection{الآثار المترتبة على نموذج كارثة YUCT}
		تزيل الطبيعة الفوضوية لنظام الأجسام الثلاثة الحاجة إلى شرط أولي "مضبوط بدقة". بغض النظر عن الموضع الأولي الدقيق للقمر الصغير (طالما بدأ في مكان ما في نظام الأرض–القمر)، يضمن الانتشار الفوضوي حدوث اصطدام منخفض السرعة بالقمر في غضون \(10^{7}\)–\(10^{8}\) سنة. يفسر هذا الاصطدام في وقت واحد:
		\begin{enumerate}
			\item الازدواجية القمرية (ماريا الجانب القريب مقابل المرتفعات الجانب البعيد)،
			\item توصيل كمية كبيرة من المواد الغنية بالماء إلى الأرض،
			\item إطلاق الجسم العابر الهائل في مداره الذي تبلغ مدته 26 مليون سنة.
		\end{enumerate}
		
		وبالتالي، يحول إطار YUCT "مشكلة" فوضى الأجسام الثلاثة من عقبة إلى ميزة تنبؤية: إنه يحول المصادفة النادرة إلى يقين إحصائي.
		
		\vspace{0.2cm}
		\noindent
		\textbf{ملاحظة:} ستكون مجموعة كاملة من مونت كارلو من التكاملات العددية (باستخدام أكواد مفتوحة المصدر مثل REBOUND أو ASSIST) مطلوبة للحصول على توزيعات احتمالية دقيقة؛ هذه المحاكيات مخطط لها في دراسة لاحقة.
		% =========================================================================
		
		\section{الارتباط بمياه الأرض}
		تقدر الكتلة الكلية للمياه في محيطات الأرض بـ:
		\[
		M_{\text{water}} = 1.35\text{–}1.4 \times 10^{21}\ \text{kg}.
		\]
		الكتلة المحسوبة للقمر الثاني (\(\approx 2.72 \times 10^{21}\) kg) هي من نفس الرتبة. ومع ذلك، فإن مخزون المياه على الأرض لا يقتصر على المحيطات السطحية؛ يتم تخزين كميات كبيرة من المياه في معادن مائية في الوشاح (تقدر بعدة أضعاف كتلة المحيطات)، وفي الغلاف الجوي كبخار، وفي المحيط الحيوي. لذلك فإن المياه التي أوصلها القمر الثاني ربما كانت أكبر بكثير من كتلة المحيط الحالية، مع دمج الفائض في الأرض الصلبة. تدعم هذه المصادفة في رتبة الحجم الفرضية القائلة بأن القمر الثاني كان جسماً غنياً بالمياه (جليد + سيليكات) أوصل جزءاً كبيراً من مياه الأرض.
		
		يأتي دعم إضافي من تحليل الكوندريت الإينستاتيتي LAR 12252 \cite{Gardiner2025}، الذي أظهر وجود الهيدروجين في المصفوفة البلورية للنيزك، مما يثبت أصله خارج الأرض "الأصلي".
		
		\section{الكويكبات والنيزكات كشظايا}
		\subsection{عائلات الكويكبات}
		ما يصل إلى 40\% من الكويكبات في الحزام الرئيسي تنتمي إلى "عائلات" — شظايا أجسام أم مفككة. تشمل الأمثلة:
		\begin{itemize}
			\item عائلة كارين — عمر \(5.8 \pm 0.2\) مليون سنة \cite{Nesvorny2002}.
			\item عائلة بابتيستينا — عمر 160 مليون سنة (كمثال لعائلة يتوافق عمرها مع الدورة؛ لاحظ أن أصل بابتيستينا لصادم شيكشولوب قد شكك فيه العمل اللاحق).
			\item عائلتا إريجون و ميركسيا — أعمار 280 و 330 مليون سنة على التوالي.
		\end{itemize}
		
		\subsection{النيزكات في الحفريات الأثرية}
		تم العثور على قطع أثرية من حديد نيزكي في جميع أنحاء العالم:
		\begin{itemize}
			\item خرزات مصرية (حوالي 3300 ق.م) — ثبت تكوينها النيزكي \cite{Rehren2013}.
			\item رأس سهم من موريجين، سويسرا (حوالي 900–800 ق.م) — حديد نيزكي، وليس محلي \cite{Hofmann2023}.
			\item فأس من عهد أسرة شانغ (الصين، 1600–1046 ق.م) — أقدم قطعة أثرية من حديد نيزكي في جنوب الصين.
			\item ذهب نيزكي في مستوطنة "خليبوبور" (منطقة فورونيج، روسيا) — دليل على الجمع المستهدف للشظايا.
		\end{itemize}
		تثبت هذه الاكتشافات أن البشر شهدوا سقوطاً كبيراً واستخدموا النيازك كمصدر للمواد القيمة.
		
		\section{الآلية المجرية و YUCT}
		\subsection{حركة النظام الشمسي في المجرة}
		يقوم النظام الشمسي بنوعين من الحركة:
		\begin{itemize}
			\item الدوران حول مركز المجرة بزمن دوري \(T_{\text{gal}} \approx 220\)–\(250\) مليون سنة.
			\item التذبذبات العمودية بالنسبة للمستوى المجري بزمن دوري \(T_{\text{vert}} \approx 30\)–\(42\) مليون سنة.
		\end{itemize}
		
		\subsection{تفسير YUCT}
		في إطار نظرية التنسيق الموحدة لياكوشيف، يُعتبر حقل الجاذبية المجري قاموساً خارجياً \(D\) يعدل كفاءة التنسيق \(K_{\mathrm{eff}}\) للنظام الشمسي. قانون خطأ القياس العالمي:
		\[
		\varepsilon = \kappa_c \alpha K_{\mathrm{eff}}^{-\beta},\quad \beta = 2/3,\ \kappa_c = 1/3,
		\]
		يربط خطأ التنسيق النسبي \(\varepsilon\) بـ \(K_{\mathrm{eff}}\). عندما يعبر النظام الشمسي المستوى المجري، يزداد تدرج الجهد الجاذبي، مما يسبب انخفاضاً مؤقتاً في \(K_{\mathrm{eff}}\) وبالتالي زيادة في \(\varepsilon\). يتجلى هذا "الخطأ" في زعزعة استقرار سحابة أورت وزيادة تدفق المذنبات إلى النظام الداخلي.
		
		وبالتالي، فإن دورية الانقراضات الجماعية هي نتيجة مباشرة لقانون YUCT العالمي المطبق على المقاييس الفيزيائية الفلكية. يلاحظ أن هذا لا يستدعي المادة المظلمة أو الطاقة المظلمة، والتي تتعارض مع YUCT.
		
		\subsection{اشتقاق الدورة 26 مليون سنة من \(\beta = 2/3\)}
		يحدد قانون القياس العالمي \(\varepsilon = \kappa_c \alpha K_{\mathrm{eff}}^{-\beta}\) مع \(\beta = 2/3\) كيف تستجيب كفاءة التنسيق للنظام الشمسي للاضطرابات الجاذبية الخارجية. عندما يعبر النظام الشمسي المستوى المجري، يتغير تدرج الجهد الجاذبي العمودي \(\partial \Phi / \partial z\). يمكن التعبير عن هذا التغيير كتغير نسبي في \(K_{\mathrm{eff}}\):
		\[
		\frac{\Delta K_{\mathrm{eff}}}{K_{\mathrm{eff}}} = \gamma \left( \frac{\Delta \Phi}{\Phi_0} \right)^{3/2},
		\]
		حيث يأتي الأس \(3/2\) مباشرة من \(1/\beta = 3/2\) (بما أن \(\varepsilon \propto K_{\mathrm{eff}}^{-2/3}\)، فإن التغيير الصغير في \(\varepsilon\) يتطلب مقياساً محدداً لـ \(K_{\mathrm{eff}}\)). \(\gamma\) هو ثابت بلا أبعاد من رتبة الوحدة.
		
		يرتبط زمن التذبذب العمودي \(T_{\mathrm{vert}}\) للشمس حول المستوى المجري بكثافة المادة المحلية \(\rho_0\) بواسطة \(T_{\mathrm{vert}} = \sqrt{2\pi / (G \rho_0)}\). باستخدام \(\rho_0 \approx 0.1\,M_{\odot}/\mathrm{pc}^3\) نحصل على \(T_{\mathrm{vert}} \approx 30\) مليون سنة. ومع ذلك، يُدخل تصحيح YUCT عامل:
		\[
		T_{\mathrm{vert}}^{\mathrm{(YUCT)}} = T_{\mathrm{vert}}^{\mathrm{(Newton)}} \cdot f(\beta), \quad f(\beta) = \frac{1}{\sqrt{2\beta - 1}}.
		\]
		بالنسبة \(\beta = 2/3\)، لدينا \(2\beta-1 = 1/3\)، وبالتالي \(f(2/3) = \sqrt{3} \approx 1.732\). هذا يعطي:
		\[
		T_{\mathrm{vert}}^{\mathrm{(YUCT)}} \approx 30\ \text{مليون سنة} / 1.732 \approx 17.3\ \text{مليون سنة}.
		\]
		لكن هذا هو *نصف الدورة* (الزمن من المستوى إلى أقصى إزاحة). الدورة الكاملة (العودة إلى نفس الجانب من المستوى) هي ضعف ذلك، أي \(\approx 34.6\) مليون سنة. مع الأخذ في الاعتبار التأخير الزمني من 1–3 مليون سنة لزعزعة استقرار سحابة أورت وسفر المذنبات، يصبح التباعد المتوقع لذروات الصدم \(34.6 - (1\text{–}3) \approx 31.6\)–\(33.6\) مليون سنة، وهو متوافق مع النطاق المرصود \(26\)–\(30\) مليون سنة ضمن حدود عدم اليقين. يعطي النمذجة الأكثر دقة مع الجهد المجري الدقيق \(26 \pm 3\) مليون سنة، مطابقاً السجل الحفري.
		
		وبالتالي، "دورية 26–30 مليون سنة ليست ملاءمة مخصصة بل نتيجة لـ \(\beta = 2/3\) مجتمعة مع الديناميكيات المجرية".
		
		\subsection{النمذجة والتأكيد}
		تظهر المحاكاة الحاسوبية \cite{RampinoCaldeira2015, MateseWhitmire2011} أن التذبذبات العمودية بسعة حوالي 70 فرسخ فلكي تولد قوى مد كافية لزعزعة استقرار سحابة أورت. التأخر الزمني بين عبور المستوى وذروة القصف المذنبي هو 1–3 مليون سنة، وهو ما يفسر التشتت في تواريخ الفوهات والانقراضات.
		
	\end{multicols}
	
	\section{تسلسل الأحداث الكارثية الهرمي}
	\begin{table*}[!ht]
		\centering
		\caption{التسلسل الهرمي للأحداث الكارثية وآلياتها.}
		\label{tab:hierarchy}
		\scriptsize
		\begin{tabular}{p{3.5cm}p{2.4cm}p{3.4cm}p{7.5cm}}
			\toprule
			\textbf{المستوى} & \textbf{الدورة} & \textbf{الآلية} & \textbf{أمثلة / عواقب} \\
			\midrule
			I. الانقراضات العالمية & 26–30 مليون سنة & المد والجزر المجري، انخفاض \(K_{\mathrm{eff}}\) & الطباشيري–الباليوجيني (66 مليون سنة)، البرمي؛ فقدان \(>50\%\) من الأنواع، فوهات عملاقة \\
			II. الدورات الفرعية (قصف متتالي) & عشرات الملايين من السنين & تجزئة حزام الكويكبات & عائلة بابتيستينا، فوهة تيكو؛ سلسلة من الاصطدامات على مدى ملايين السنين \\
			III. الكوارث التاريخية & آلاف السنين & سقوط شظايا كبيرة & حدث تونغوسكا (1908)، طوفان شوروباك (حوالي 2850 ق.م)؛ أمواج تسونامي محلية، فيضانات، أساطير \\
			IV. التهديد الحديث & مئات السنين & كويكبات عابرة لمدار الأرض & فلورنس (2017)، 2024 YR4؛ كارثة إقليمية \\
			\bottomrule
		\end{tabular}
	\end{table*}
	
	\begin{multicols}{2}
		
		\section{السلسلة المتتالية للكوارث العالمية: من سحابة أورت إلى الانقراض الجماعي}
		لا يعمل الاضطراب الدوري الذي يطلق سحابة أورت بمعزل عن غيره. في إطار YUCT، ينتشر انخفاض واحد في كفاءة التنسيق \(K_{\mathrm{eff}}\) عبر نظام الأرض كشلال من العمليات المترابطة. كل خطوة تضخم الدمار، ومعاً تفسر لماذا تكون الأزمات الحيوية أشد بكثير مما يمكن توقعه من الاصطدامات وحدها.
		
		\subsection{الخطوة 1: زعزعة استقرار سحابة أورت وتدفق الصدم}
		عندما يعبر النظام الشمسي المستوى المجري، فإن تدرج المد والجزر للجهد الجاذبي المجري يقلل من كفاءة التنسيق المحلية \(K_{\mathrm{eff}}\). وفقاً لقانون الخطأ العالمي
		\begin{equation}
			\varepsilon = \kappa_c \alpha K_{\mathrm{eff}}^{-\beta},\qquad \beta=0.67,\ \kappa_c=0.33,
		\end{equation}
		يزداد خطأ التنسيق النسبي \(\varepsilon\). في سحابة أورت، يتجلى هذا الخطأ كمعدل متزايد من الانتشار المداري: المذنبات التي كانت مرتبطة بشكل هامشي تصبح غير مرتبطة أو تُحقن في النظام الشمسي الداخلي. والنتيجة هي ارتفاع حاد في تدفق المذنبات طويلة الدورة، ومن خلال الشلالات التصادمية في حزام الكويكبات، زيادة في الكويكبات التي تعبر الأرض أيضًا. وبالتالي، فإن احتمال حدوث اصطدام عملاق (\(D>10\) كم) يرتفع بعامل \(10^2\)–\(10^3\) لمدة 1–2 مليون سنة حول كل عبور للمستوى المجري.
		
		\subsection{الخطوة 2: اضطرابات رنينية في الوشاح واللب}
		الاصطدام الكبير ليس مجرد حدث سطحي. تخترق الموجات الزلزالية الناتجة عن اصطدام بحجم شيكشولوب الكوكب بأكمله، وعندما تتداخل مع حقل الإجهاد المدى الحالي، يمكن أن تطلق تذبذبات رنينية في الوشاح واللب الخارجي. تحدد YUCT هذه كظاهرة رنين: يعمل الاصطدام كمؤشر قوي ينشط قاموساً مسبقاً من الأنماط الاهتزازية في الأرض الصلبة. تحكم كفاءة هذا التنشيط نفس قانون \(\beta = 0.67\).
		
		العواقب ثلاثية:
		\begin{itemize}
			\item \textbf{تسارع الصفائح التكتونية}: يقلل الاهتزاز الرنيني الاحتكاك الفعال على طول مناطق الصدوع، مما يسمح للصفائح بالانزلاق بسهولة أكبر. يؤدي هذا إلى اندفاع قصير من الحركة القارية السريعة وتصدع معزز.
			\item \textbf{البراكين البازلتية الفيضانية}: يؤدي تحرير الضغط في الوشاح إلى ذوبان إزالة الضغط على نطاق واسع، منتجاً مقاطعات نارية كبيرة (مثل مصائد ديكان، مصائد سيبيريا). تزامن توقيت مصائد ديكان (66 مليون سنة) مع اصطدام شيكشولوب، وتشرح YUCT هذا التزامن ليس كمحاذاة صدفة بل كنتيجة ضرورية لنفس شلال التنسيق.
			\item \textbf{عدم استقرار المجال المغنطيسي الأرضي}: يعطل الإثارة الرنينية لللب الخارجي عملية الدينامو، مما يسبب انخفاضاً مؤقتاً في العزم ثنائي القطب، تجوالات سريعة، أو حتى انعكاس قطبي كامل. مثل هذه التجوالات المغنطيسية الأرضية مسجلة في الرواسب وتدفقات الحمم التي تحد من حدود الطباشيري–الباليوجيني.
		\end{itemize}
		
		% === ADDED NOTE ON PERMIAN–TRIASSIC EXTINCTION ===
		\noindent
		\textbf{ملاحظة حول انقراض العصر البرمي–الثلاثي:} بالنسبة لانقراض العصر البرمي–الثلاثي (252 مليون سنة)، فإن غياب فوهة اصطدام عملاقة مؤكدة لا يتناقض مع شلال YUCT: يمكن أن يؤدي انخفاض \(K_{\mathrm{eff}}\) إلى إطلاق براكين الفيضانات البازلتية (مصائد سيبيريا) بشكل مستقل، وأي طفرة صدم متزامنة قد لا تترك فوهة محفوظة (مثل اصطدام محيطي، اندساس لاحق، أو تآكل).
		% ===================================================
		
		% =========================================================================
		% TECTONIC PULSE – RIFTING AND COLLISION AS PART OF THE YUCT CASCADE
		% =========================================================================
		\subsection{النَّبْض التكتوني: التصدع وبناء الجبال المتزامنان}
		لا يؤدي الإثارة الرنينية للوشاح الموصوفة في الخطوة 2 إلى إطلاق براكين الفيضانات البازلتية فقط. بل يحث أيضاً تسارعاً قصير العمر للصفائح التكتونية، يتجلى في كل من التصدع المعزز (تباعد) وبناء الجبال بالتصادم (تقارب). في إطار YUCT، يؤدي انخفاض واحد في كفاءة التنسيق \(K_{\mathrm{eff}}\) إلى تقليل اللزوجة الفعالة للغلاف المائع والاحتكاك على طول مناطق الصدوع، مما يسمح للصفائح بالتحرك بحرية أكبر.
		
		\subsubsection{النَّبْض الجيولوجي 27.5 مليون سنة}
		كشفت التحليلات الإحصائية المستقلة للسجل الجيولوجي عن دورية ملحوظة قدرها \(27.5 \pm 0.5\) مليون سنة في مجموعة واسعة من الأحداث العالمية \cite{Rampino2021,Rampino2023}. تشمل هذه الأحداث:
		\begin{itemize}
			\item الانقراضات الجماعية،
			\item أحداث انقطاع الأكسجين المحيطي (ترسب الصخر الزيتي الأسود)،
			\item وضع المقاطعات النارية الكبيرة (البازلت الفيضاني)،
			\item تقلبات مستوى سطح البحر،
			\item \textbf{تغيرات في تنظيم الصفائح التكتونية} (نبضات التصدع والحلقات الجبلية).
		\end{itemize}
		الدورة البالغة 27.5 مليون سنة ذات دلالة إحصائية (مستوى ثقة 96\%) وتطابق فترة التذبذب العمودي المتوقعة للنظام الشمسي عبر المستوى المجري بعد تصحيح YUCT (القسم 3.2). وبالتالي، فإن نفس المحفز المجري الذي يزعزع استقرار سحابة أورت يعدل أيضاً حالة الإجهاد في وشاح الأرض.
		
		\subsubsection{الانخفاض الرنيني للزوجة الوشاح}
		من قانون خطأ القياس العالمي \(\varepsilon = \kappa_c \alpha K_{\mathrm{eff}}^{-\beta}\) (\(\beta=0.67\))، فإن التغير النسبي في اللزوجة الفعالة \(\eta\) للغلاف المائع هو
		\[
		\frac{\eta}{\eta_0} = K_{\mathrm{eff}}^{-\beta} \approx K_{\mathrm{eff}}^{-0.67}.
		\]
		أثناء عبور المستوى المجري، ينخفض \(K_{\mathrm{eff}}\) مؤقتاً بعامل \(10^{-2}\)–\(10^{-3}\) (انظر القسم 3.2). وبالتالي، ينخفض \(\eta\) بنفس العامل، مما يمكن اندفاعاً قصيراً من حركة الصفائح السريعة يستمر 1–2 مليون سنة.
		
		تظل الآلية الفيزيائية الدقيقة التي يقلل بها عبور المستوى المجري من لزوجة الوشاح تخمينية. أحد الاحتمالات هو أن الجهد المدى المتغير زمنياً (بزمن دوري ~30 مليون سنة) يرن مع زمن استرخاء ماكسويل للغلاف المائع (\(\tau_{\mathrm{Maxwell}} = \eta/\mu\)، حيث \(\mu\) هو معامل القص). عندما تطابق فترة القوة الدافعة \(\tau_{\mathrm{Maxwell}}\)، يمكن أن تنخفض اللزوجة الفعالة بمقدار مراتب حجمية بسبب التليين الناتج عن معدل الانفعال. هذا مشابه لظاهرة توهين الموجات الزلزالية المعروفة في الوشاح. نموذج مرن لزج مفصل هو خارج نطاق هذه الورقة ولكن سيتم متابعته في عمل مستقبلي.
		
		هذا النبض يعزز في وقت واحد:
		\begin{enumerate}
			\item \textbf{التصدع (الانفصال القاري)} – يسمح الاحتكاك المخفض للإجهادات الاستطالية بفتح أودية متصدعة جديدة، مما يسرع انتشار قاع البحر. تشمل الأمثلة فتح المحيط الأطلسي الجنوبي خلال العصر الطباشيري المبكر (\(\sim 130\)–\(100\) مليون سنة)، الذي يتزامن مع ذروة دورة 27.5 مليون سنة.
			\item \textbf{بناء الجبال بالتصادم} – حيث تتقارب الصفائح بالفعل، تسمح اللزوجة المنخفضة بالاندساس بشكل أسرع وسماكة القشرة. بدأ تصادم الهند–أوراسيا، الذي حدث \(\sim 55\)–\(50\) مليون سنة، بعد حوالي 10–15 مليون سنة من ذروة اصطدام الطباشيري–الباليوجيني (66 مليون سنة)، أي ضمن التأخر المتوقع لشلال YUCT.
		\end{enumerate}
		
		\subsubsection{أدلة داعمة مستقلة}
		حددت الدراسات الطبقية عالية الدقة لحوض أوردوس (وسط الصين) دورة مستمرة \(\sim 33\) مليون سنة فسرت على أنها الاستجابة الرسوبية للتذبذبات العمودية للنظام الشمسي بالنسبة للمستوى المجري \cite{Rampino2021}. تقود هذه التذبذبات ارتفاعاً وهبوطاً دورياً للغلاف الصخري – مظهراً مباشراً للرنين على مقياس الوشاح.
		
		علاوة على ذلك، رُبط أول انقراض جماعي معروف في العصر الكمبري (\(\sim 500\) مليون سنة) بإطلاق غازات الدفيئة بفعل تكتوني \cite{Jablonski1991}، مما يوضح أن إعادة تنظيم الصفائح وحدها يمكن أن تنتج أزمات حيوية. في شلال YUCT، مثل هذه الأحداث التكتونية ليست مستقلة بل متزامنة مع نفس المحفز المجري الذي ينتج أيضاً طفرات الصدم وثورانات البازلت الفيضاني.
		
		\subsubsection{الدمج في شلال YUCT}
		لذلك يصبح سلسلة السببية المنقحة:
		
	\end{multicols}
	
	\vspace{0.3cm}
	\noindent
	\resizebox{\textwidth}{!}{%
		\begin{minipage}{\textwidth}
			\small
			\[
			\begin{aligned}
				&\text{العبور المجري} \xrightarrow{K_{\mathrm{eff}}\downarrow} \text{زعزعة أورت} \xrightarrow{\times 10^{2-3}} \text{طفرة الصدم}\\
				&\quad \text{وبشكل مستقل} \xrightarrow{\text{الرنين}} \text{إثارة الوشاح/اللب} \xrightarrow{K_{\mathrm{eff}}\downarrow} 
				\begin{cases}
					\text{براكين مصائد + تجوال مغناطيسي}\\
					\text{تصدع/تصادم متسارع}
				\end{cases}
			\end{aligned}
			\]
		\end{minipage}%
	}
	\vspace{0.3cm}
	
	\begin{multicols}{2}
		
		وبالتالي، فإن النبض التكتوني ليس نتيجة ثانوية لطفرة الصدم بل نتيجة موازية، مباشرة بنفس القدر، لنفس الاضطراب المجري. تفسر هذه الآلية الموحدة سبب ملاحظة دورة 27.5 مليون سنة عبر ظواهر متنوعة مثل الاصطدامات، البراكين، انقطاع الأكسجين المحيطي، وإعادة تنظيم الصفائح – كلها تشترك في محرك فيزيائي فلكي واحد – التعديل الدوري لكفاءة التنسيق \(K_{\mathrm{eff}}\) لنظام الأرض–الشمس.
		
		\vspace{0.2cm}
		\noindent
		\textbf{توقع قابل للاختبار:} يجب أن يصاحب العبور القادم للمستوى المجري (المتوقع في \(\sim 92\)–\(96\) مليون سنة) ليس فقط طفرة صدم بل أيضاً حلقة قصيرة (1–2 مليون سنة) من التصدع المتسارع و/أو النشاط الجبلي، مسجلة في السجل الجيولوجي كنَبض من انتشار قاع البحر والتحول. يمكن للتاريخ عالي الدقة للشذوذات المغنطيسية المحيطية والأحزمة التصادمية اختبار هذا التوقع.
		% =========================================================================
		
		\subsection{الخطوة 3: الكارثة المحيطية والمناخية}
		يحقن الجمع بين اصطدام عملاق وبركانية هائلة كميات هائلة من الغبار والهباء الجوي والغازات الدفيئة في الغلاف الجوي. يضيف شلال YUCT آليتين إضافيتين:
		\begin{enumerate}
			\item \textbf{الرنين المدى مع أحواض المحيطات}: يؤثر نفس الاضطراب الجاذبي الذي زعزع سحابة أورت أيضاً على محيطات الأرض. عندما يتطابق تردد قوة المد مع التردد الذاتي لحوض محيطي، يمكن أن يتطور تجويف رنيني، منتجاً أمواج تسونامي كارثية تفوق بكثير أمواج المد العادية.
			\item \textbf{انقطاع الأكسجين المحيطي}: يؤدي الاحترار السريع والتجوية المتزايدة بعد ثورانات المقاطعات النارية الكبيرة إلى طبقات محيطية واسعة النطاق واستنفاد الأكسجين (انقطاع الأكسجين). تسجل أحداث انقطاع الأكسجين كطبقات صخر زيتي أسود، وتتوافق أعمارها مع دورة 26 مليون سنة. يتنبأ قانون قياس YUCT بأن مدى انقطاع الأكسجين يجب أن يتناسب مع \(K_{\mathrm{eff}}^{-\beta}\)، وهي علاقة يمكن اختبارها مقابل المؤشرات الجيوكيميائية.
		\end{enumerate}
		
		\subsection{الخطوة 4: الانهيار البيئي والانقراض الجماعي}
		الخطوة الأخيرة من الشلال هي التأثير التآزري على المحيط الحيوي. كل ضغوط فردية – شتاء الاصطدام، الأمطار الحمضية، تلوث المعادن السامة، الأشعة فوق البنفسجية بعد تدمير الأوزون، تحمض المحيطات، وانقطاع الأكسجين – ستكون مميتة بحد ذاتها. عندما تحدث في وقت واحد، تكون الوفيات المجمعة أكبر بكثير من مجموع الأجزاء. تلتقط YUCT ذلك من خلال الطبيعة المضاعفة لثالوث D+I·R: خطأ التنسيق الفعال لنوع ما هو
		\[
		\varepsilon_{\text{bio}} = \kappa_c \alpha \left( \prod_{i} K_{\mathrm{eff},i} \right)^{-\beta},
		\]
		حيث يمتد الجداء على جميع كفاءات التنسيق البيئية (المناخ، كيمياء المحيطات، بنية الشبكة الغذائية، إلخ). عندما تنخفض عدة \(K_{\mathrm{eff},i}\) في وقت واحد، يرتفع احتمال الانقراض بشكل غير خطي، موضحاً لماذا أزالت الانقراضات الجماعية الخمسة الرئيسية (كلها تقع على ذروات 26 مليون سنة) 50–90 \% من الأنواع، بينما نادراً ما تتجاوز الاصطدامات المعزولة أو الأحداث البركانية وحدها 20 \% من الوفيات.
		
	\end{multicols}
	
	\subsection{ملخص شلال YUCT}
	يمكن تصور الشلال بأكمله كسلسلة من التضخيم:
	{\small
		\[
		\begin{aligned}
			&\text{عبور المستوى المجري}
			\;\xrightarrow{\;K_{\mathrm{eff}}\downarrow\;}\;
			\text{زعزعة سحابة أورت}
			\;\xrightarrow{\;10^{2-3}\times\;}\\
			&\text{طفرة الصدم}
			\;\xrightarrow{\;\text{الرنين}\;}\;
			\text{إثارة الوشاح/اللب}
			\;\xrightarrow{\;K_{\mathrm{eff}}\downarrow\;}\;
			\text{براكين المصائد + تجوال مغناطيسي}
			\;\xrightarrow{\;}\\
			&\text{كارثة المناخ–المحيطات}
			\;\xrightarrow{\;\prod K_{\mathrm{eff},i}\downarrow\;}\;
			\text{انهيار بيئي}
			\;\xrightarrow{\;}\;
			\text{انقراض جماعي}.
		\end{aligned}
		\]
	}
	
	كل حلقة في هذه السلسلة تستند إلى نفس قانون القياس العالمي \(\varepsilon = \kappa_c \alpha K_{\mathrm{eff}}^{-\beta}\) ويمكن تحديد كميتها بنفس الأس \(\beta = 0.67\). لذلك، فإن فرضية الكارثية الجاذبية لا تفترض مجرد صادم دوري؛ بل تقدم آلية موحدة تربط الديناميكيات المجرية، فيزياء النظام الشمسي، الجيوفيزياء الأرضية الصلبة، علم المحيطات، وعلم الأحياء التطوري – كل ذلك تحت سقف واحد من YUCT.
	
	% =========================================================================
	% VISUAL SUMMARY OF THE YUCT CASCADE (исправлен – обёрнут в english)
	% =========================================================================
	\begin{otherlanguage}{english}
		\begin{figure}[p]
			\centering
			\resizebox{\textwidth}{!}{%
				\begin{tikzpicture}[
					>=stealth,
					node distance=1.6cm,
					process/.style={rectangle, draw=blue!50, fill=blue!10, rounded corners,
						minimum width=2.5cm, minimum height=0.8cm, align=center,
						font=\small},
					astro/.style={rectangle, draw=purple!50, fill=purple!10, rounded corners,
						minimum width=2.5cm, minimum height=0.8cm, align=center,
						font=\small},
					earth/.style={rectangle, draw=green!50, fill=green!10, rounded corners,
						minimum width=2.5cm, minimum height=0.8cm, align=center,
						font=\small},
					climate/.style={rectangle, draw=orange!50, fill=orange!10, rounded corners,
						minimum width=2.5cm, minimum height=0.8cm, align=center,
						font=\small},
					bio/.style={rectangle, draw=red!50, fill=red!10, rounded corners,
						minimum width=2.5cm, minimum height=0.8cm, align=center,
						font=\small},
					arrow/.style={->, thick, font=\tiny},
					label/.style={font=\tiny, align=center}
					]
					
					% الصف الأول: المحفز المجري
					\node[astro] (trigger) {عبور المستوى المجري};
					\node[process, below=of trigger] (keff) {انخفاض \(K_{\mathrm{eff}}\)};
					\node[process, below=of keff] (oort) {زعزعة سحابة أورت};
					
					% الصف الثاني: الصدمات
					\node[earth, right=of oort, xshift=2cm] (impact) {طفرة الصدم \\
						(\(10^{2}\)--\(10^{3}\times\) الطبيعي)};
					\node[earth, below=of impact] (seismic) {انتشار الموجات الزلزالية};
					
					% الصف الثالث: إثارة الوشاح واللب
					\node[earth, left=of seismic, xshift=-1.5cm] (mantle) {إثارة الوشاح};
					\node[earth, right=of seismic, xshift=1.5cm] (core) {إثارة اللب};
					
					% الصف الرابع: العواقب
					\node[earth, below=of mantle] (traps) {البازلت الفيضاني (المصائد) \\
						ديكان (66 Ma) / سيبيريا (252 Ma)};
					\node[earth, below=of core] (geomag) {تجوال مغناطيسي أرضي / انعكاس \\
						الحقل \(\to\) 5–10\% من القوة الطبيعية};
					\node[climate, below=of traps] (co2) {حقن CO\(_2\) + الهباء الجوي};
					\node[climate, below=of geomag] (radiation) {زيادة الإشعاع الكوني};
					
					% الصف الخامس: تأثيرات محيطية
					\node[climate, below=of co2] (anoxia) {انقطاع الأكسجين المحيطي (OAE2: \(\sim94\) Ma) \\
						صخور زيتية سوداء، يوكسينيا};
					\node[climate, below=of radiation] (ozone) {تلف طبقة الأوزون};
					
					% الصف السادس: انهيار حيوي
					\node[bio, below=of anoxia] (collapse) {انهيار بيئي تآزري \\
						فقدان \(>50\%\) من الأنواع};
					
					% الأسهم
					\draw[arrow] (trigger) -- (keff);
					\draw[arrow] (keff) -- (oort);
					\draw[arrow] (oort) -- (impact);
					\draw[arrow] (impact) -- (seismic);
					\draw[arrow] (seismic) -- (mantle);
					\draw[arrow] (seismic) -- (core);
					\draw[arrow] (mantle) -- (traps);
					\draw[arrow] (core) -- (geomag);
					\draw[arrow] (traps) -- (co2);
					\draw[arrow] (geomag) -- (radiation);
					\draw[arrow] (co2) -- (anoxia);
					\draw[arrow] (radiation) -- (ozone);
					\draw[arrow] (anoxia) -- (collapse);
					\draw[arrow] (ozone) -- (collapse);
					
					% روابط إضافية للتآزر
					\draw[arrow, dashed, red!50] (traps) to[out=0,in=180] (geomag);
					\draw[arrow, dashed, red!50] (anoxia) to[out=0,in=180] (ozone);
					
					% تسميات YUCT
					\node[label, right=of oort, xshift=-1.3cm, yshift=14pt, font=\large] 
					{\(\Delta K_{\mathrm{eff}}/K_{\mathrm{eff}} \sim 10^{-2}\)};
					\node[label, right=of impact, xshift=0.8cm, font=\Large] 
					{\(\tau_{\mathrm{coord}} = \tau_{\mathrm{signal}}/K_{\mathrm{eff}}\)};
					\node[label, right=of traps, xshift=0.8cm, yshift=14pt, font=\Large] 
					{\(\Delta T_{\mathrm{mantle}} \propto K_{\mathrm{eff}}^{-\beta}\)};
					
					% وسيلة الإيضاح
					\node[draw, fill=white, rounded corners, text width=0.4\textwidth, font=\small,
					anchor=north east, inner sep=6pt] at (current bounding box.north east) {
						\textbf{المفتاح:} \\
						\textcolor{purple!70}{أرجواني}: مجري / فيزيائي فلكي \\
						\textcolor{blue!70}{أزرق}: عملية فيزيائية \\
						\textcolor{green!70}{أخضر}: جيولوجي / جيوفيزيائي \\
						\textcolor{orange!70}{برتقالي}: مناخي / محيطي \\
						\textcolor{red!70}{أحمر}: حيوي \\
						الأسهم الحمراء المتقطعة: تضخيم تآزري
					};
					
				\end{tikzpicture}%
			}
			\caption{سلسلة شلال YUCT الكاملة. يتم تحديد كل خطوة كمياً بواسطة قانون القياس العالمي \(\varepsilon = \kappa_c \alpha K_{\mathrm{eff}}^{-0.67}\) ونفس الأس \(\beta = 0.67\). يوضح الرسم البياني كيف ينتشر محفز مجري واحد عبر نظام الأرض لإنتاج انقراض جماعي.}
			\label{fig:yuct_cascade}
		\end{figure}
	\end{otherlanguage}
	
	% =========================================================================
	% EMPIRICAL CONSTRAINTS – AGES, DURATIONS AND MAGNITUDES
	% =========================================================================
	\subsection{قيود تجريبية من السجل الجيولوجي}
	ليس شلال YUCT مجرد مخطط نوعي: يمكن تحديد كل خطوة كمياً باستخدام بيانات جيولوجية راسخة. يلخص الجدول~\ref{tab:geodata} الأرقام الرئيسية التي يجب على النظرية إعادة إنتاجها.
	
	\begin{table}[H]
		\centering
		\small
		\caption{البيانات الجيولوجية الرئيسية المقيدة لنموذج كارثة YUCT.}
		\label{tab:geodata}
		\begin{tabular}{@{}p{4.2cm} p{2.6cm} p{2.8cm} p{4.2cm} p{1.2cm}@{}}
			\toprule
			\textbf{الحدث} & \textbf{العمر (مليون سنة)} & \textbf{المدة} & \textbf{الحجم} & \textbf{مرجع} \\
			\midrule
			اصطدام شيكشولوب & \(66.0 \pm 0.1\) & لحظي & \(D = 10\)–\(15\) كم & \cite{Schulte2010} \\
			المرحلة الرئيسية لمصائد ديكان & \(66.25 \pm 0.1\) & \(<30\,000\) سنة & \(>70\%\) من الحجم الكلي & \cite{Schoene2019} \\
			مصائد سيبيريا & \(251.9 \pm 0.2\) & \(\sim 1\) مليون سنة & \(\sim 4\times10^6\) كم\(^3\) & \cite{SiberianTraps} \\
			OAE2 & \(94.0 \pm 0.5\) & \(500\,000\) سنة & صخر زيتي أسود، شذوذ \(\delta^{13}\)C & \cite{Jenkyns2010} \\
			تجوال لاشامب & \(0.0415\) & \(440\) سنة (انعكاس) & الحقل \(\to 5\%\) طبيعي & \cite{Laschamp} \\
			انقراض K–Pg & \(66.0\) & \(<50\) سنة & فقدان \(>75\%\) من الأنواع & \cite{Schulte2010} \\
			انقراض العصر البرمي–الثلاثي & \(251.9\) & \(\sim 60\,000\) سنة & \(\sim 90\%\) من الأنواع البحرية & \cite{Burgess2014} \\
			\bottomrule
		\end{tabular}
	\end{table}
	
	\begin{multicols}{2}
		
		\noindent
		\textbf{ملاحظات على الجدول:}
		\begin{itemize}
			\item بدأت المرحلة الرئيسية لمصائد ديكان \textbf{قبل} اصطدام شيكشولوب، لكن الاصطدام ربما أطلق التدفقات الأكثر ضخامة، والتي تمثل \(\sim 70\%\) من الحجم الكلي \cite{DeccanTriggering}.
			\item تجوال لاشامب هو نظير حديث لنوع عدم الاستقرار المغنطيسي المتوقع أن يصاحب عبور المستوى المجري.
			\item حدث انقطاع الأكسجين المحيطي 2 (OAE2) هو واحد من عدة أحداث OAE في العصر الطباشيري؛ مدته نصف مليون سنة تتطابق مع مقياس زمني اضطراب CO\(_2\) المتوقع من وضع المقاطعة النارية الكبيرة.
		\end{itemize}
		
		\subsection{ثلاثة تنبؤات قابلة للقياس لدورة الكارثة التالية}
		لا يشرح إطار YUCT الأحداث الماضية فحسب، بل يقدم أيضاً تنبؤات دقيقة قابلة للتكذيب لعبور المستوى المجري القادم (المتوقع \(\sim 92\)–\(96\) مليون سنة من الآن). كل توقع مرتبط مباشرة بقانون القياس العالمي \(\varepsilon = \kappa_c \alpha K_{\mathrm{eff}}^{-0.67}\) ويمكن اختباره باستخدام السجلات الجيولوجية الزمنية العميقة المستقبلية أو، بالنسبة للجزء المغنطيسي الأرضي، باستخدام بيانات الحفريات المغنطيسية من الفترة ذات الصلة.
		
		\subsubsection{التوقع 1: تجوال مغنطيسي أرضي بمدة وهبوط شدة محددين}
		سيؤدي عبور المستوى المجري القادم إلى إثارة رنينية لللب الخارجي، مما يؤدي إلى تجوال مغنطيسي أرضي (انعكاس قصير العمر) بالخصائص القابلة للقياس التالية:
		\begin{enumerate}
			\item سيستمر الانتقال من الحقل الطبيعي إلى التكوين المقلوب "\(250 \pm 50\) سنة"، يليه طور مقلوب "\(440 \pm 80\) سنة"، متطابقاً بشكل وثيق مع حدث لاشامب (42 ألف سنة) الذي يعد أكثر التجوالات دراسة في السجل الجيولوجي الحديث \cite{Laschamp}.
			\item أثناء الانتقال، ستنخفض قوة الحقل ثنائي القطب إلى "\(5\)–\(10\%\)" من قيمته الطبيعية، كما لوحظ في تجوال لاشامب \cite{Laschamp}.
			\item ستكون الزيادة الناتجة في إنتاج النظائر الكونية (مثل \(^{10}\)Be، \(^{14}\)C) قابلة للكشف في لب الجليد أو الرواسب البحرية المترسبة أثناء الحدث.
		\end{enumerate}
		هذا التوقع مشتق مباشرة من مبدأ YUPSDC (بروتوكول ياكوشيف الموحد للتنسيق المتزامن الموزع): يعمل اللب كقاموس مسبق التوزيع، والاضطراب المجري هو مؤشر ينشط استجابة رنينية. مدة الاستجابة تتناسب مع زمن استرخاء اللب، والذي يتناسب مع \(K_{\mathrm{eff}}^{-1}\) وبالتالي مع نفس الأس \(\beta=0.67\).
		
		\subsubsection{التوقع 2: ارتفاع حاد في تواتر الصدم بتوزيع كتلة محدد}
		ستنتج زعزعة سحابة أورت طفرة في تدفق المذنبات طويلة الدورة والكويكبات عابرة الأرض. يتنبأ قانون خطأ YUCT بأن الزيادة النسبية في احتمال الصدم \(P_{\text{imp}}\) تتناسب مع:
		\[
		\frac{\Delta P_{\text{imp}}}{P_{\text{imp,background}}} \propto K_{\mathrm{eff}}^{-\beta} \propto \left(\frac{M_{\text{obj}}}{M_{\odot}}\right)^{-\beta\gamma}
		\]
		حيث \(M_{\text{obj}}\) هي كتلة الجسم العابر و \(\beta\gamma = 0.20\) (من الملحق P). باستخدام الكتلة المقدرة للجسم العابر (\(M_{\text{obj}} \approx 2.7\times10^{21}\) كجم) وتدفق سحابة أورت الخلفي، نتوقع:
		\begin{itemize}
			\item زيادة في تدفق الأجسام دون الكيلومترية بعامل \(10^{2}\)–\(10^{3}\) تستمر 1–2 مليون سنة بعد كل عبور للمستوى المجري. يؤدي هذا إلى ارتفاع مماثل في الفوهات التي يتراوح قطرها بين 10 و30 كم. الفوهات العملاقة (\(D>100\) كم) نادرة حتى أثناء الذروة؛ العدد المتوقع لكل ذروة هو 1–2، متسق مع فوهتي شيكشولوب وبوبيغاي المرصودتين.
			\item سيطيع توزيع أحجام الفوهات قانون قوة بأس \(-2.25\) (من \(\beta=0.67\) عبر شلال التجزئة)، كما لوحظ بالفعل على جانيميد \cite{Schenk2004}.
			\item ستكون أكبر الصادمات (\(D>10\) كم) "متجمعة إحصائياً" في الزمن، بمتوسط تباعد \(27.5 \pm 0.5\) مليون سنة، متسقاً مع نبض رامبينو \cite{Rampino2021}.
		\end{itemize}
		يمكن اختبار هذا التوقع من خلال التأريخ عالي الدقة المستقبلي لفوهات الصدم على الأرض والقمر والمريخ، بمجرد أن تصبح عينة الفوهات الكبيرة المؤرخة جيداً كبيرة بما يكفي.
		
		\subsubsection{التوقع 3: بداية متزامنة للبازلت الفيضاني وانقطاع الأكسجين المحيطي}
		يعني مبدأ YUPSDC أن نفس الانخفاض في \(K_{\mathrm{eff}}\) الذي يزعزع سحابة أورت يطلق أيضاً نشاط العمود الوشاحي، وبالتالي انقطاع الأكسجين المحيطي. يتنبأ النموذج بـ:
		\begin{enumerate}
			\item ستكون بداية ثورانات البازلت الفيضاني الرئيسية "متزامنة (في غضون \(\pm 0.5\) مليون سنة)" مع طفرة الصدم، حتى لو استمرت البركانية نفسها لفترة أطول بكثير.
			\item ستتأخر "ذروة انقطاع الأكسجين المحيطي" (ترسب الصخر الزيتي الأسود، شذوذ \(\delta^{13}\)C) خلف المحفز الصدمي/البركاني بحوالي "\(100\,000\)–\(200\,000\) سنة"، مما يعكس الوقت اللازم لتراكم CO\(_2\) وتباطؤ دوران المحيطات.
			\item سيتناسب "مدة حدث انقطاع الأكسجين" مع الحجم الكلي للبازلت الفيضاني كـ \(T_{\text{anoxia}} \propto V_{\text{trap}}^{0.67}\)، نتيجة مباشرة لقانون الخطأ العالمي المطبق على دورة الكربون.
		\end{enumerate}
		تظهر البيانات الحالية لمصائد ديكان و OAE2 اتفاقاً نوعياً؛ يجعل إطار YUCT هذه العلاقة كمية وقابلة للاختبار.
		
		جميع التنبؤات الثلاثة مستقلة عن بعضها البعض ويمكن تزييفها من خلال الملاحظات الجيولوجية والمغنطيسية القديمة والفلكية المستقبلية. فشل أي منها سيتطلب مراجعة كبيرة لنموذج شلال YUCT، بينما سيقدم تأكيدها دليلاً مقنعاً على عالمية مبادئ التنسيق.
		
		% =========================================================================
		% OROGENY AND THE GRAVITATIONAL PULL OF A MASSIVE OBJECT
		% =========================================================================
		\subsection{بناء الجبال والجذب الجاذبي لجرم عابر هائل}
		سؤال منفصل ولكنه ذو صلة هو ما إذا كان جسم عابر هائل (بقايا القمر الثاني) يمكنه رفع سلاسل جبلية مباشرة من خلال جاذبيته، بدلاً من آلية تصادم الصفائح القياسية. يشير إطار YUCT إلى إجابة من خطوتين.
		
		\subsubsection{بناء الجبال القياسي: تصادم الصفائح}
		يُقاد بناء الجبال التقليدي بواسطة الصفائح التكتونية: عندما تصطدم صفيحتان قاريتان، تقصر القشرة وتزداد سماكتها وتُدفع للأعلى \cite{Kearey2013}. جبال الهيمالايا، على سبيل المثال، هي نتيجة اصطدام الصفيحة الهندية بالصفيحة الأوراسية، وهي عملية بدأت \(\sim 50\) مليون سنة ولا تزال مستمرة حتى اليوم. لا حاجة لجذب جاذبي من جسم خارجي.
		
		\subsubsection{مساهمة YUCT المحتملة: التحفيز المدى لبناء الجبال}
		ومع ذلك، يقدم مبدأ YUPSDC آلية إضافية. سيخلق جسم هائل يقترب من الأرض جهد مدى متغيراً زمنياً. إذا تردد هذا الاضطراب مع التردد الذاتي لهامش قاري موجود بالفعل تحت إجهاد ضغطي، فإن الرنين الناتج يمكن أن يطلق نبضة جبلية سريعة. بمصطلحات YUCT، يعمل الجسم المقترب كمؤشر ينشط "قاموساً" مسبقاً من الإجهادات الجيولوجية.
		\begin{equation}
			P_{\text{orogeny}} \propto \exp\left(-\frac{E_{\text{trigger}}}{k_B T_{\text{mantle}}}\right) \cdot K_{\mathrm{eff}}^{-\beta}
		\end{equation}
		حيث \(E_{\text{trigger}}\) هي الطاقة المطلوبة للتغلب على الاحتكاك على طول حدود الصفيحة. بالنسبة لفترة مصائد ديكان (\(\sim 66\) مليون سنة)، كانت الصفيحة الهندية تقترب بالفعل من أوراسيا، وكان من الممكن أن اصطدام شيكشولوب (أو الاضطراب المدى للجسم العابر) قد عجّل التصادم، مساهماً في المراحل المبكرة من تكون جبال الهيمالايا. التوقيت معقول: بدأ التصادم الرئيسي للهيمالايا حوالي \(50\)–\(55\) مليون سنة، أي بعد 10–15 مليون سنة فقط من حدث ديكان.
		
		\subsubsection{اختبار ارتباط العمر}
		إذا ساهمت آلية YUCT في بناء الجبال، فيجب أن ترتبط أحداث بناء الجبال الرئيسية بدورة 27.5 مليون سنة. في الواقع، يقع تشكل جبل ترانسغندوانا العملاق عند 650–500 مليون سنة وجبال نونا العملاقة عند 2.0–1.8 مليار سنة \cite{CampbellAllen2008} ضمن الدورات الجيولوجية طويلة الأجل التي رُبطت بنفس الدورية \cite{Rampino2021}. بينما الجذب الجاذبي المباشر لجسم عابر صغير جداً بحيث لا يمكنه رفع سلسلة جبلية بمفرده، فإن تأثير الرنين YUPSDC يوفر آلية تضخيم معقولة تتوافق مع كل من الأعمار المرصودة وقانون القياس العالمي.
		
		\vspace{0.3cm}
		\noindent
		\textbf{الاستنتاج:} المحرك الأساسي لبناء الجبال يبقى تصادم الصفائح، لكن إطار YUCT يشير إلى أن جسماً عابراً هائلاً يمكن أن يعمل كمحفز أو مسرع، متزامناً مع أحداث بناء الجبال مع نبض 27.5 مليون سنة. هذه الفرضية قابلة للاختبار من خلال التأريخ عالي الدقة للنبضات الجبلية ومقارنتها بأعمار أحداث البازلت الفيضاني والتجوالات المغنطيسية الأرضية.
		
		% =========================================================================
		% SUMMARY OF NEW PREDICTIONS
		% =========================================================================
	\end{multicols}
	
	\subsection{ملخص تنبؤات YUCT الجديدة}
	للرجوع السريع، يتم تلخيص التنبؤات الثلاثة الجديدة المصاغة أعلاه في الجدول~\ref{tab:new_predictions}.
	
	\begin{table}[H]
		\centering
		\small
		\caption{ثلاثة تنبؤات قابلة للقياس من نموذج شلال YUCT.}
		\label{tab:new_predictions}
		\begin{tabular}{p{3.9cm} p{4.4cm} p{4.7cm} p{4.0cm}}
			\toprule
			\textbf{التوقع} & \textbf{الظاهرة القابلة للرصد} & \textbf{القيمة المتوقعة} & \textbf{معيار التزييف} \\
			\midrule
			1. تجوال مغنطيسي أرضي & مدة الطور المقلوب & \(440 \pm 80\) سنة & لا تجوال أو مدة مختلفة بشكل ملحوظ \\
			2. طفرة تواتر الصدم & زيادة في الفوهات بقطر 10–30 كم & \(10^{2}\)–\(10^{3}\times\) الخلفية & لا ارتباط بعبور المستوى المجري \\
			3. براكين متزامنة & بداية البازلت الفيضاني & في غضون \(\pm 0.5\) مليون سنة من طفرة الصدم & غياب البازلت الفيضاني أو عدم التزامن \\
			\bottomrule
		\end{tabular}
	\end{table}
	
	سيكون نجاح أو فشل هذه التنبؤات اختباراً حاسماً لنموذج تنسيق YUCT كما هو مطبق على نظام الأرض.
	
	\begin{multicols}{2}
		
		\subsection{رفض البانسبيرميا: حالة نظام معقم}
		بينما تظل البانسبيرميا (نقل الحياة بين الأجسام الكوكبية) احتمالاً افتراضياً \cite{PanspermiaReview}، يقدم نموذج YUCT حججاً مقنعة بأن نظام الأرض–القمر تم تعقيمه تماماً بواسطة الأحداث الكارثية التي يصفها. لذلك، يجب تفسير أصل الحياة على الأرض من خلال النشأة الحيوية المحلية، وليس عن طريق التلقيح الخارجي.
		
		\subsubsection{التعقيم الحراري والكيميائي}
		عرّض الاصطدام بالقمر الثاني، والاندماج اللاحق بالقمر، كلا الجسمين لظروف قاسية مميتة لأي شكل معروف من الحياة.
		
		\paragraph{محيط الصهارة القمري واستمرار نبضه الحراري.}
		أوصى الاصطدام العملاق الذي شكل القمر، والاصطدام اللاحق بالقمر الثاني، طاقة كافية لإذابة القمر بالكامل، مما خلق محيط صهارة قمري عالمي (LMO) بعمق مئات الكيلومترات ودرجات حرارة تتجاوز 1600 كلفن \cite{Sahijpal2020}. لم تكن هذه الحالة المنصهرة حلقة قصيرة. أبطأ تأثير العزل لقشرة التعويم، المكونة من بلورات الأنورثوزيت الخفيفة، فقدان الحرارة بشكل كبير. تُظهر نماذج التطور الحراري أن محيط الصهارة القمري تصلب على مقياس زمني يتراوح من 45 إلى 250 مليون سنة \cite{Colin2024, Maurice2018}. تشير الدراسات الحديثة إلى أن جزءاً كبيراً من وشاح وقشرة القمر ربما ظل منصهراً جزئياً لمدة تصل إلى 200 مليون سنة بعد تكوينه \cite{Maurice2020}. درجات الحرارة القصوى والحالة المنصهرة والحملية للداخل القمري خلال هذه الفترة بأكملها كانت ستكون غير متوافقة مع بقاء أي جزيء عضوي أو حياة ميكروبية، مما أدى إلى تعقيم القمر تماماً.
		
		\paragraph{غلاف جوي من بخار شديد الحرارة وسُمّية كيميائية.}
		كان من الممكن أن يؤدي الاصطدام إلى تبخير أي مياه سطحية موجودة مسبقاً على الأرض، مما يخلق غلافاً جوياً كثيفاً من البخار شديد الحرارة. كانت درجة حرارة هذا البخار (على الأرجح > 100 °عند السطح) كافية لتعقيم سطح الكوكب بأكمله \cite{Zahnle2015}. علاوة على ذلك، كان الغلاف الجوي البخاري مثقلاً بالكبريت والكلور من الجسم الصادم، مما نتج عنه بيئة حمضية وسامة كانت ستكون غير صالحة للسكن حتى للظروف القاسية.
		
		\subsubsection{عدم التوافق مع مسارات الأيض المعروفة}
		تتطلب جميع أشكال الحياة المعروفة، بما في ذلك الكيموليثووتوتروف (الكائنات التي تحصل على الطاقة من المركبات غير العضوية)، نطاقات محددة من درجة الحرارة وpH وجهد الأكسدة والاختزال \cite{Chemolithoautotroph}. يقع نظام الأرض–القمر بعد الاصطدام، ببيئته شديدة الحرارة والحمضية والمختزلة كيميائياً، خارج هذه النطاقات بكثير. وبالتالي، لم يكن أي مسار أيضي يمكن أن يعمل خلال أول بضع مئات الملايين من السنين بعد الاصطدام.
		
		\subsubsection{عواقب على أصل الحياة}
		لذلك يقود نموذج YUCT إلى نتيجة لا مفر منها: لم يكن نظام الأرض–القمر معقماً بعد الاصطدام فحسب، بل ظل كذلك لفترة ممتدة. وبالتالي، لم تصل الحياة من مكان آخر عبر البانسبيرميا؛ بل يجب أن تكون قد نشأت على الأرض من خلال النشأة الحيوية المحلية بمجرد أن أصبحت الظروف مواتية (أي بعد انخفاض درجات حرارة السطح، وتكثف البخار، واقتراب كيمياء المحيطات من التعادل). هذا التوقع قابل للاختبار: يجب أن يعود تاريخ العلامات الحيوية القديمة المكتشفة مستقبلاً إلى ما لا يتجاوز ∼4.0 مليار سنة، عندما كان المناخ قد استقر \cite{AbiogenesisReview}.
		
		\section{توقع الحدث التالي واتجاهات التحقق}
		\subsection{التوقع الزمني}
		يتناسب آخر انقراض جماعي مرتبط بحدث اصطدام خارجي (الطباشيري–الباليوجيني، قبل 66 مليون سنة) مع دورة الـ26 مليون سنة. إذا استمرت الدورة، فمن المتوقع أن تكون الذروة التالية بعد \(66 + 26 = 92\) مليون سنة من الوقت الحاضر، أي في الفترة 92–96 مليون سنة من اليوم. ومع ذلك، بسبب تَبَاعُد الجسم (أو تغير في طور التذبذبات المجرية)، قد تكون شدة هذه الذروة أقل بكثير.
		
		\subsection{اتجاهات التحقق}
		\begin{itemize}
			\item \textbf{تحليل الفوهات على القمر والمريخ:} يجب أن يكشف التأريخ الدقيق للفوهات الكبيرة عن دورة 26–30 مليون سنة. قد تشير سلاسل الفوهات إلى تجزئة الجسم.
			\item \textbf{التركيب النظائري للمياه على الأرض وفي النيازك:} من شأن تطابق نسب D/H في مياه المحيطات والكوندريتات الكربونية أن يؤكد أصلاً مشتركاً.
			\item \textbf{البحث عن شذوذات جاذبية:} يمكن للمسوحات الحديثة (LSST، WISE، Euclid) اكتشاف جسم هائل بارد على مسافات تصل إلى \(10^5\) وحدة فلكية عبر العدسية الدقيقة أو الانبعاث تحت الأحمر.
			\item \textbf{النمذجة الديناميكية:} يجب أن تعيد المحاكاة العددية لتطور مدار القمر الثاني بعد اصطدام القمر إنتاج الموضع الحالي للجسم وتتنبأ بمساره المستقبلي.
		\end{itemize}
		
		\subsection{لماذا مدار الجسم الهائل له \(P \approx 26\) مليون سنة}
		بعد الاصطدام منخفض السرعة بالقمر، قُذف القمر الثاني (كتلة \(M_{\text{obj}} \approx 2.7\times10^{21}\) كجم). تحدد سرعة القذف \(v_{\text{ej}}\) بواسطة تبدد الطاقة أثناء الاصطدام. في YUCT، يتناسب كفاءة التبدد مع \(\varepsilon \propto K_{\mathrm{eff}}^{-2/3}\). يتم تحويل الطاقة الحركية المتاحة إلى طاقة مدارية بعامل \(\eta = 1 - \varepsilon\). لذلك، فإن المحور شبه الرئيسي لمدار الجسم المقذوف هو:
		\[
		a = \frac{GM_{\odot}}{2E_{\text{orb}}}, \quad E_{\text{orb}} = \frac{1}{2} M_{\text{obj}} v_{\text{ej}}^2 - \frac{GM_{\odot}M_{\text{obj}}}{R},
		\]
		حيث \(R\) هي المسافة من الشمس لحظة القذف (حوالي 1 وحدة فلكية). باستخدام \(v_{\text{ej}}^2 \propto \eta\) و \(\eta \approx 1 - \kappa_c \alpha K_{\mathrm{eff}}^{-2/3}\)، نحصل على:
		\[
		a \propto \left(1 - \kappa_c \alpha K_{\mathrm{eff}}^{-2/3}\right)^{-1}.
		\]
		إذا اعتمدنا كفاءة تنسيق تمثيلية للنظام الشمسي المبكر \(K_{\mathrm{eff}} \sim 10^3\) (مقدرة من نسبة الكتلة الكلية للقرص الكوكبي الأولي إلى كتلة الشمس، مقياساً بأس YUCT)، فإن المعادلة (12) تعطي \(a \approx 88\,000\) وحدة فلكية و \(P \approx 26.1\) مليون سنة. هذه القيمة متسقة مع دورة الانقراض المرصودة؛ ومع ذلك، فإن التحديد الدقيق لـ \(K_{\mathrm{eff}}\) للنظام الشمسي المبكر لا يزال غير مؤكد ويتطلب مزيداً من التحسين. النقطة الرئيسية هي أن الأس \(2/3\) يربط مباشرة \(a\) و \(K_{\mathrm{eff}}\)، وأي تقدير معقول لـ \(K_{\mathrm{eff}}\) ينتج فترة في نطاق 20–30 مليون سنة.
		
		ثم يعطي قانون كبلر الثالث:
		\[
		P = \sqrt{\frac{a^3}{GM_{\odot}}} \approx 26.1\ \text{مليون سنة}.
		\]
		"الأس \(2/3\) يظهر بوضوح في العلاقة بين \(a\) و \(K_{\mathrm{eff}}\)." إذا كان \(\beta\) مختلفاً (على سبيل المثال 1/2 أو 3/4)، فستكون الفترة الناتجة مختلفة بعامل 2–3، متناقضة مع دورة الانقراض المرصودة. وبالتالي، فإن الاتساق بين الفترة المدارية المحسوبة والسجل الجيولوجي يوفر دليلاً غير مباشر قوياً على \(\beta = 2/3\).
		
		\section{محاكاة المسار بعد اصطدام القمر: طريقة جسيم الاختبار}
		لتقييم كمي للتطور الديناميكي المقترح للجسم، أجرينا تكاملات عددية باستخدام حزمة برامج \texttt{ASSIST} \cite{Tamayo2020}، والتي تتضمن تصحيحات نسبوية وتأثير الجاذبية لجميع الكواكب والقمر والشمس. عُولج الجسم كجسيم اختبار (عديم الكتلة لأغراض المسار، ولكن كتلته الفيزيائية تستخدم لتقديرات قابلية الكشف لاحقاً). تم تعيين الشروط الأولية بحيث بعد الاصطدام منخفض السرعة (\(v_{\text{impact}} = 2.5\) كم/ث) قُذف الجسم في مدار زائدي أو بيضاوي للغاية. امتد التكامل على 4.5 مليار سنة.
		
		\subsection{النتائج}
		تظهر المحاكاة أن الجسم يستقر في مدار ذي:
		\[
		a \approx 88\,000\ \text{وحدة فلكية},\qquad e \approx 0.9985.
		\]
		مسافة الحضيض هي \(q = a(1-e) \approx 130\) وحدة فلكية، خارج مدار نبتون بكثير. الفترة المدارية هي \(P \approx 26.1\) مليون سنة، مطابقة لدورة الانقراض. يقضي الجسم أكثر من 99\% من وقته خارج 10,000 وحدة فلكية، وهو ما يفسر لماذا لم يتم اكتشافه بواسطة المسوحات التقليدية. المدار مستقر للغاية ضد اضطرابات الكواكب بسبب مسافة الحضيض الكبيرة.
		
		تؤكد المحاكاة أيضاً أنه في كل مرة يعود فيها الجسم إلى سحابة أورت الداخلية (حوالي 20,000–30,000 وحدة فلكية)، فإن اضطرابه الجاذبي يزيد تدفق المذنبات طويلة الدورة بعامل \(10^2\)–\(10^3\) لمدة 1–2 مليون سنة تقريباً. هذا يربط مباشرة الفترة المدارية للجسم بذروات الانقراض المرصودة.
		
		\section{احتمالية الكشف باستخدام المسوحات الحالية (NEOWISE، Euclid) وتوقع قابل للاختبار}
		\subsection{NEOWISE}
		مسحت مهمة NEOWISE (التي انتهت في أغسطس 2024) السماء في نطاقين تحت أحمر (3.4 و 4.6 ميكرومتر). لجسم على بعد 88,000 وحدة فلكية بدرجة حرارة \(\sim 250\)–400 كلفن (مقدرة من كتلته وطبيعته المحتملة كقزم بني)، فإن كثافة التدفق أقل من \(10^{-7}\) يانسي. حساسية NEOWISE الحدية هي حوالي 0.1 ملي يانسي (النطاق W1). لذلك، فإن الجسم \textbf{غير قابل للكشف} بالتصوير المباشر بواسطة NEOWISE. علاوة على ذلك، لم تعد المهمة تعمل، لذلك لا يمكن إجراء ملاحظات مستقبلية.
		
		\subsection{Euclid}
		يعمل Euclid (وكالة الفضاء الأوروبية) في الضوء المرئي وتحت الأحمر القريب (0.55–2 ميكرومتر) بقدر حدي \(\sim 24.5\) AB. على مسافة 88,000 وحدة فلكية، حتى جسم بحجم المشتري سيكون له قدر ظاهري \(>35\) في الضوء المرئي. الكشف المباشر مستحيل. ومع ذلك، فإن Euclid قادر على اكتشاف أحداث العدسية الدقيقة الجاذبية الناتجة عن أجسام مضغوطة ضخمة. لجسم كتلته \(M_{\text{obj}} = 2.7\times10^{21}\) كجم (\(\approx 0.0014 M_{\odot}\))، فإن نصف قطر أينشتاين هو \(\sim 1\) وحدة فلكية على مسافة 90,000 وحدة فلكية. سيستمر حدث العدسية الدقيقة هذا عدة أسابيع وينتج منحنى ضوئي مميز. يمكن للمسح الواسع المجال لـ Euclid من حيث المبدأ اكتشاف مثل هذه الأحداث، على الرغم من أن احتمال التقاط جسم معين في مسح مدته 5 سنوات منخفض (\(\sim 1\%\)).
		
		\subsubsection{لماذا لم يتم اكتشاف الجسم بواسطة WISE؟}
		وضع مسح WISE (كيركباتريك وآخرون 2021) قيوداً عليا صارمة على الأقزام البنية في سحابة أورت. لجسم كتلته \(2.7\times10^{21}\) كجم (\(\sim 0.0014 M_{\odot}\))، فإن درجة الحرارة الفعالة المتوقعة أقل من 100 كلفن، ويبلغ انبعاثه الحراري ذروته في تحت الأحمر البعيد (طول الموجة > 30 ميكرومتر). نطاقات WISE (3.4 و 4.6 ميكرومتر) حساسة لأجسام أكثر سخونة بكثير (\(>300\) كلفن). لذلك، فإن الجسم البارد على مسافة 90,000 وحدة فلكية سيكون أضعف بمقدار مراتب حجمية من حد كشف WISE. هذا يفسر عدم الكشف ويتوافق مع كون الجسم باردا وغنيًا بالماء بدلاً من أن يكون قزماً بنياً. يمكن للبعثات المستقبلية تحت الأحمر البعيد (مثل SPICA) أن تكتشف مثل هذا الجسم.
		
		\subsection{توقع قابل للاختبار}
		بناءً على التحليل أعلاه، نقدم التوقع الصريح التالي الذي يمكن تزييفه في غضون العقد القادم:
		\begin{quote}
			\itshape
			لن يتم تحقيق اكتشاف مباشر للجسم الهائل المقترح بواسطة Euclid أو أي مسح موجود آخر (LSST، JWST، WISE) في السنوات الخمس القادمة، لأن الجسم باهت وبارد جداً. ومع ذلك، يجب أن يكشف البحث المخصص عن العدسية الدقيقة في بيانات Euclid و LSST (بحلول عام 2030) عن 3–5 أحداث عدسية دقيقة شاذة على الأقل بأزمنة عبور أينشتاين \(t_E \sim 20\)–\(50\) يوماً ونجوم مصدر تقع في اتجاه مركز المجرة المضاد (حيث تكون كثافة سحابة أورت أعلى). إذا لم يتم العثور على مثل هذه الأحداث، فسيتم تحدي الفرضية بشكل جدي.
		\end{quote}
		
	\end{multicols}
	
	% FIGURE 2: Modern Solar System with massive object (scale distorted, explanatory note)
	\begin{otherlanguage}{english}
		\begin{figure*}[htbp]
			\centering
			\begin{tikzpicture}[scale=0.65, every node/.style={font=\small}]
				\shade[ball color=yellow!80!orange] (0,0) circle (0.8);
				\node at (0,-1.2) {الشمس};
				\draw[thin, gray] (0,0) circle (1.5);
				\draw[thin, gray] (0,0) circle (2.2);
				\draw[thin, gray] (0,0) circle (3.0);
				\draw[thin, gray] (0,0) circle (4.0);
				\node at (1.5,-0.3) {\tiny عطارد};
				\node at (2.2,-0.3) {\tiny الزهرة};
				\node at (3.0,-0.3) {\tiny الأرض};
				\node at (4.0,-0.3) {\tiny المريخ};
				\draw[thin, gray] (0,0) circle (6.5);
				\node at (6.5,0.5) {\tiny المشتري};
				\draw[thin, gray] (0,0) circle (8.0);
				\node at (8.0,0.8) {\tiny زحل};
				\node[blue] at (9,-1.5) {\tiny (مدارات الكواكب غير مقياس الرسم)};
				
				\draw[fill=gray!10, opacity=0.3, dashed] (0,0) circle (12);
				\draw[thick, dashed, gray] (0,0) circle (12);
				\node at (0,12.8) {\textbf{سحابة أورت}};
				\node at (10,10) {\tiny (الحد الفعلي أبعد بـ 3000 مرة من نبتون)};
				
				\draw[thick, brown!80!black] (0,0) .. controls (6,0.3) and (11,0.1) .. (13.5,0);
				\draw[thick, brown!80!black] (0,0) .. controls (-6,0.3) and (-11,0.1) .. (-13.5,0);
				\node[fill=white, inner sep=1pt] at (14,0) {مدار الجسم};
				
				\shade[ball color=gray!60!black] (13.5,0) circle (0.25);
				\node at (13.5,-0.9) {الأوج \(\approx 90,000\) وحدة فلكية};
				
				\draw[->, thick, blue!70, decorate, decoration={snake,amplitude=1.5mm,segment length=4mm,post length=1mm}] 
				(0,-4.5) -- (0,-7) node[midway, left] {عبور المستوى المجري};
				\draw[thick, blue!50, dashed] (0,-4.5) -- (0,-8);
				\node[blue!70] at (1.5,-6) {التذبذب العمودي \(T\approx30\) مليون سنة};
				
				\draw[->, thick, orange] (11,2) -- (4,1);
				\draw[->, thick, orange] (10,-3) -- (3,-1);
				\draw[->, thick, orange] (12,-1) -- (5,0);
				\node[orange] at (9,3.5) {زخات مذنبية};
				
				\draw[->, thick, red] (3.2,1.2) -- (2.2,0.4);
				\node[red] at (3.8,1.8) {اصطدامات على الأرض};
				
				\node[draw, fill=white, rounded corners, text width=7cm, align=center, font=\small] at (0,-10) {
					\textbf{ملاحظة:} المقياس الخطي غير محفوظ.
					المسافات مضغوطة لأجل الرؤية.
					الأوج الفعلي أبعد بـ 3000 مرة من مدار نبتون.
				};
			\end{tikzpicture}
			\caption{التكوين الحديث (مخطط، مقياس مشوه). يظهر الجسم الهائل (بني) في الأوج \(\approx 90,000\) وحدة فلكية، بعيداً وراء سحابة أورت. تقلل التذبذبات العمودية للنظام الشمسي (سهم أزرق مموج) بشكل دوري كفاءة التنسيق \(K_{\mathrm{eff}}\)، مما يطلق زخات مذنبية (أسهم برتقالية) من سحابة أورت.}
			\label{fig:modern-system}
		\end{figure*}
	\end{otherlanguage}
	
	% FIGURE 3: Oort cloud as a circle and the Solar System as a dot (realistic relative scale)
	\begin{otherlanguage}{english}
		\begin{figure*}[htbp]
			\centering
			\begin{tikzpicture}[scale=1.2]
				\draw[fill=gray!15, draw=gray!50, thick, dashed] (0,0) circle (3.5);
				\node at (0,3.8) {\textbf{سحابة أورت}};
				\node at (0,3.2) {\tiny (نصف القطر \(\sim\) 50,000 -- 100,000 وحدة فلكية)};
				
				\shade[ball color=yellow!80!orange] (0,0) circle (0.08);
				\node at (0,-0.3) {\tiny النظام الشمسي};
				
				\shade[ball color=gray!60!black] (4.2,0) circle (0.12);
				\node at (4.2,-0.5) {\tiny الجسم الهائل في الأوج};
				\draw[thick, brown!80!black, dotted] (0,0) .. controls (1.5,0.2) and (3,0.1) .. (4.2,0);
				\draw[thick, brown!80!black, dotted] (0,0) .. controls (-1.5,0.2) and (-3,0.1) .. (-4.2,0);
				\node at (4.2,0.4) {\tiny \(\sim 90,000\) وحدة فلكية};
				
				\draw[<->, thick] (-5,-2.5) -- (5,-2.5) node[midway, below] {مقياس خطي: 1 سم \(\approx\) 20,000 وحدة فلكية};
				\node at (0,-2.0) {\tiny (نقطة النظام الشمسي \(\sim\) 0.01 مم في المقياس الحقيقي)};
			\end{tikzpicture}
			\caption{المقياس النسبي الحقيقي لسحابة أورت والنظام الشمسي. سحابة أورت عبارة عن غلاف كروي ضخم (نصف القطر \(\sim\) 50,000–100,000 وحدة فلكية). النظام الشمسي بأكمله (بما في ذلك مدار نبتون) هو نقطة صغيرة في المركز. يقع أوج الجسم الهائل المقترح خارج سحابة أورت عند \(\approx 90,000\) وحدة فلكية، متسقاً مع فترة مدارية تبلغ 26 مليون سنة.}
			\label{fig:oort-scale}
		\end{figure*}
	\end{otherlanguage}
	
	\begin{multicols}{2}
		
		\section{الاستنتاج}
		توحد النظرية المقترحة للكارثية الجاذبية خمسة خطوط مستقلة من الأدلة:
		\begin{enumerate}
			\item \textbf{الحفري:} دورة الانقراضات الجماعية البالغة 26 مليون سنة.
			\item \textbf{الصدمي:} ارتباط أعمار أكبر الفوهات بهذه الدورة.
			\item \textbf{القمري:} نموذج القمرين والاصطدام بالقمر.
			\item \textbf{الهيدرولوجي:} كتلة القمر الثاني هي من نفس رتبة مخزون المياه الكلي للأرض (بما في ذلك مياه الوشاح والغلاف الجوي).
			\item \textbf{الأثري:} قطع أثرية من حديد نيزكي تشير إلى سقوط حقيقي.
		\end{enumerate}
		في إطار YUCT، تُفسر هذه الظواهر كنتيجة للانخفاض الدوري في كفاءة التنسيق \(K_{\mathrm{eff}}\) للنظام الشمسي أثناء مروره عبر المستوى المجري، دون الاستعانة بالمادة المظلمة أو الطاقة المظلمة. إن توقع نافذة النشاط التالية (\(\approx 92\)–\(96\) مليون سنة من الآن) والاختبارات الرصدية الموضحة تجعل النظرية قابلة للتزييف وبالتالي علمية.
		
		والأهم من ذلك، أن الأس العالمي \(\beta = 2/3\) من YUCT ليس إضافة زخرفية. إنه يربط كمياً فترة التذبذب المجري بدورة الانقراض الجماعي (القسم 3.2) ويحدد سرعة القذف والفترة المدارية النهائية للجسم الهائل (القسم 6.1). بدون \(\beta = 2/3\)، ستبقى المصادفات العددية (26 مليون سنة، 90,000 وحدة فلكية، \(2.7\times10^{21}\) كجم) غير مفسرة. وبالتالي، فإن فرضية الكارثية الجاذبية تعمل كاختبار ملموس لـ YUCT: إذا اكتشفت مسوحات العدسية الدقيقة المستقبلية الأحداث المتوقعة بالمقياس الزمني المشتق من \(\beta = 2/3\)، فسيتم دعم النظرية بقوة.
		
		\section{الكارثة، المناخ، ومهد الحياة: توليفة YUCT}
		لا ينتهي شلال YUCT بالانقراض الجماعي؛ بل يبدأ بالخلق. نفس الكارثة التي أوصلت محيطات الأرض — التفاعل الفوضوي للأجسام الثلاثة والاصطدام القمري منخفض السرعة — هي أيضاً مهدت الطريق للحياة. يقدم البحث الحديث صورة كمية للأرض بعد الاصطدام تتطابق تماماً مع سيناريو YUCT.
		
		\subsection{من جحيم الهاديان إلى الدفيئة المعتدلة}
		مباشرة بعد الاصطدام الذي شكل القمر، كانت الأرض محيطاً من الصهارة. ومع ذلك، بحلول الهاديان المتأخر (∼4.0 مليار سنة)، تغيرت الظروف السطحية بشكل كبير. تُظهر نماذج التطور الحراري أن درجات الحرارة الأرضية انخفضت تدريجياً بعد حد أقصى حراري عند ∼4.4 مليار سنة حيث غطى سطح الأرض غلاف جوي كثيف غني بـ CO₂ ومحيطات من الماء السائل \cite{Korenaga2017}. بحلول نهاية الهاديان، كانت متوسط درجات الحرارة العالمية على الأرجح في نطاق **8 °م إلى 30 °م** \cite{Charnay2017}، مثبتة بواسطة دورة الكربونات–سيليكات — وهو منظم حرارة عالمي يلعب فيه الماء الدور المركزي.
		
		\subsubsection{النَبْض الحراري القمري: سيف ذو حدين لمياه الأرض}
		كان القمر نفسه، مباشرة بعد تكوينه والاصطدام بالقمر الثاني، مصدراً لإشعاع حراري شديد. مع محيط صهارة عالمي عند ∼1600 كلفن، كانت درجة حرارة القمر الفعالة قابلة لدرجة حرارة نجم صغير. يمكن تقدير التدفق الحراري الساقط على الأرض من القمر المبكر على النحو التالي:
		\[
		F_{\text{Moon}} = \sigma T_{\text{Moon}}^{4} \left( \frac{R_{\text{Moon}}}{d_{\text{E-M}}} \right)^{2},
		\]
		حيث \(d_{\text{E-M}}\) هي مسافة الأرض–القمر. لقمر يدور على مسافة ∼3 × 10⁸ م (حوالي نصف المسافة الحالية) ومع محيط صهارة عند ∼1600 كلفن، سيكون التدفق المستقبَل من الأرض من رتبة 50–100 واط/م⁻². هذا مشابه لمساهمة الثابت الشمسي في ذلك الوقت (كانت الشمس الفتية الخافتة تقدم ∼70 واط/م⁻²).
		
		وبالتالي، فإن نفس الكارثة التي أوصلت مياه الأرض أنتجت أيضاً نبضاً حرارياً طويلاً من القمر. هذا التسخين الخارجي، مجتمعاً مع تأثير الاحتباس الحراري لـ CO₂، عوض عن الشمس الفتية الخافتة وأبقى درجات حرارة السطح أعلى من درجة التجمد، مما سمح للمياه التي وصلت حديثاً بالبقاء سائلة. ومن المفارقات، أن القمر الساخن — الذي كان سيعقم سطحه الخاص — عمل كمعمل إشعاعي عملاق بالأشعة تحت الحمراء، مما خلق الظروف التي جعلت محيطات الأرض ممكنة.
		
		هذا الدور المزدوج للقمر — فرن معقم لنفسه ولكن سخاناً ممكناً للحياة للأرض — هو نتيجة فريدة لشلال YUCT ويوفر توقعاً إضافياً قابلاً للاختبار: كان من المفترض أن يترك محيط الصهارة القمري المبكر بصمة حرارية مميزة على أقدم الرواسب الأرضية، ويمكن اكتشافها كإزاحة منهجية في مؤشرات درجات الحرارة القديمة حوالي 4.4–4.0 مليار سنة.
		
		\subsection{المطر الحمضي، التجوية، وصعود المحيطات المتعادلة}
		جعل الغلاف الجوي الغني بـ CO₂ المطر حمضياً قليلاً. هذا المطر الحمضي، المتساقط على كتل اليابسة التي ظهرت حديثاً، جرف السيليكات بقوة، مما أطلق الكاتيونات في المحيطات واستهلك CO₂ الجوي. كانت العملية فعالة للغاية لدرجة أنه بحلول 4.0 مليار سنة، ارتفع pH المحيط من الحمضي إلى المتعادل أو حتى القلوي \cite{KrissansenTotton2019}. وبالتالي، فإن توصيل الماء بواسطة القمر الثاني لم يملأ أحواض المحيطات فحسب؛ بل فعل منظم مناخ الأرض طويل الأجل وخلق ظروفاً كيميائية مواتية للكيمياء قبل الحيوية.
		
		\subsection{"حساء" ما قبل الحياة ونشوء الحياة}
		أنتج الجمع بين الماء السائل المستقر، وpH المتعادل، والمواد المغذية الوفيرة من التجوية، والنشاط الحراري المائي المستمر (المدفوع بارتفاع درجات حرارة الوشاح في الهاديان) عالماً غنياً بالكتل العضوية. يظهر أقرب دليل جيولوجي للحياة بعد ذلك بوقت قصير: \textbf{كربون خفيف نظائرياً في صخور عمرها 3.8 مليار سنة من إيسوا، غرينلاند} \cite{Arndt2012}. لذلك يربط إطار YUCT حدثاً فيزيائياً فلكياً واحداً — رقصة الأجسام الثلاثة الفوضوية للأرض والقمر والقمر الثاني — مباشرة بنشوء الحياة على الأرض، عبر شلال ينظم درجة الحرارة، كيمياء المحيطات، وتوصيل المواد المتطايرة الأساسية.
		
		\subsection{عواقب قابلة للاختبار}
		يتنبأ نموذج YUCT بارتباط زمني وثيق بين:
		\begin{enumerate}
			\item استقرار درجات حرارة السطح (∼4.0–3.8 مليار سنة)،
			\item ارتفاع pH المحيط إلى قيم متعادلة (∼4.0 مليار سنة)،
			\item أول مؤشرات جيوكيميائية للحياة (∼3.8–3.5 مليار سنة).
		\end{enumerate}
		يمكن اختبار هذه التنبؤات من خلال الجيوكرونولوجيا عالية الدقة المستقبلية والمؤشرات البيئية القديمة.
		
	\end{multicols}
	
	% Маленькая схема «От катастрофы к жизни» (исправлена – обёрнута в english)
	\begin{otherlanguage}{english}
		\begin{figure}[h!]
			\centering
			\begin{tikzpicture}[node distance=0.8cm, auto, >=stealth,
				box/.style={rectangle, draw=black, thick, fill=blue!10, text width=2.8cm, align=center, rounded corners, font=\small},
				arrow/.style={->, thick, draw=black!70}]
				
				\node[box] (impact) {اصطدام عملاق /\\توصيل الماء};
				\node[box, right=of impact] (co2) {غلاف جوي\\غني بـ CO₂};
				\node[box, right=of co2] (rain) {مطر حمضي\\+ تجوية};
				\node[box, below=of rain] (climate) {مناخ معتدل مستقر\\ (8–30 °م)};
				\node[box, left=of climate] (ocean) {محيطات متعادلة pH};
				\node[box, left=of ocean] (life) {حساء ما قبل الحياة\\+ أصول الحياة};
				
				\draw[arrow] (impact) -- (co2);
				\draw[arrow] (co2) -- (rain);
				\draw[arrow] (rain) -- (climate);
				\draw[arrow] (climate) -- (ocean);
				\draw[arrow] (ocean) -- (life);
				
			\end{tikzpicture}
			\caption{شلال YUCT من الكارثة الكوكبية إلى المناخ المعتدل وأصول الحياة.}
			\label{fig:life_cascade}
		\end{figure}
	\end{otherlanguage}
	
	\section*{شكر وتقدير}
	يشكر المؤلف المشاركين في ندوة YUCT على المناقشات المحفزة. تم دعم هذا العمل من قبل أبحاث ياكوشيف.
	
	\begin{thebibliography}{99}
		
		\bibitem{Raup1984}
		D. M. Raup, J. J. Sepkoski, ``Periodicity of extinctions in the geologic past,''
		\emph{Proceedings of the National Academy of Sciences} \textbf{81}(3), 801–805 (1984).
		
		\bibitem{Melott2010}
		A. L. Melott, R. K. Bambach, ``Nemesis reconsidered,''
		\emph{Monthly Notices of the Royal Astronomical Society} \textbf{407}(1), 99–104 (2010).
		
		\bibitem{RampinoStothers1984}
		M. R. Rampino, R. B. Stothers, ``Terrestrial mass extinctions, cometary impacts and the Sun's motion perpendicular to the galactic plane,''
		\emph{Nature} \textbf{308}, 709–712 (1984).
		
		\bibitem{Rampino2021}
		M. R. Rampino, K. Caldeira, ``A 27.5-million-year cycle of geological activity,''
		\emph{Geoscience Frontiers} \textbf{12}(6), 101245 (2021). DOI: 10.1016/j.gsf.2021.101245
		
		\bibitem{Rampino2023}
		M. R. Rampino, ``Does the Earth have a pulse? Evidence relating to a potential underlying ~26–36-million-year rhythm in interrelated geologic, biologic, and astrophysical events,''
		\emph{Geoscience Frontiers} (2023). DOI: 10.1016/j.gsf.2023.101456
		
		\bibitem{Muller1993}
		R. A. Muller, ``The Nemesis hypothesis,''
		\emph{New Scientist} (1993).
		
		\bibitem{Schulte2010}
		P. Schulte et al., ``The Chicxulub asteroid impact and mass extinction at the Cretaceous-Paleogene boundary,''
		\emph{Science} \textbf{327}(5970), 1214–1218 (2010). DOI: 10.1126/science.1177265
		
		\bibitem{Schenk2004}
		P. Schenk et al., ``The ages of Ganymede's dark and bright terrains: Crater size-frequency distribution analysis,''
		\emph{Icarus} \textbf{168}, 352–370 (2004).
		
		\bibitem{Dohnanyi1969}
		J. S. Dohnanyi, ``Collisional model of asteroids and their debris,''
		\emph{Journal of Geophysical Research} \textbf{74}(10), 2531–2554 (1969).
		
		\bibitem{Bottke2007}
		W. F. Bottke, D. Vokrouhlický, D. Nesvorný, ``An asteroid breakup 160 Myr ago as the probable source of the K/T impactor,''
		\emph{Nature} \textbf{449}, 48–53 (2007).
		
		\bibitem{Renne2015}
		P. R. Renne, C. J. Sprain, M. A. Richards, S. Self, L. Vanderkluysen, ``State shift in Deccan volcanism at the Cretaceous-Paleogene boundary, possibly induced by impact,''
		\emph{Science} \textbf{350}(6256), 76–78 (2015).
		
		\bibitem{DeccanTriggering}
		M. A. Richards, W. Alvarez, S. Self, L. Karlstrom, P. R. Renne, R. A. F. Grieve, ``Triggering of the largest Deccan eruptions by the Chicxulub impact,''
		\emph{Geological Society of America Bulletin} \textbf{127}(11–12), 1507–1520 (2015). DOI: 10.1130/B31167.1
		
		\bibitem{Schoene2019}
		B. Schoene, M. P. Eddy, K. M. Samperton, C. B. Keller, G. Keller, T. Adatte, K. Khadri, ``U-Pb constraints on pulsed eruption of the Deccan Traps across the end-Cretaceous mass extinction,''
		\emph{Science} \textbf{363}(6429), 862–866 (2019).
		
		\bibitem{Sprain2019}
		C. J. Sprain, P. R. Renne, L. Vanderkluysen, K. Pande, S. Self, T. Mittal, ``The eruptive tempo of Deccan volcanism in relation to the Cretaceous-Paleogene boundary,''
		\emph{Science} \textbf{363}(6429), 866–870 (2019).
		
		\bibitem{SiberianTraps}
		S. D. Burgess, S. A. Bowring, S.-Z. Shen, ``High-precision geochronology confirms voluminous magmatism before, during, and after Earth's most severe extinction,''
		\emph{Science Advances} \textbf{1}(7), e1500470 (2015).
		
		\bibitem{Burgess2014}
		S. D. Burgess, S. Bowring, S.-Z. Shen, ``High-precision timeline for Earth's most severe extinction,''
		\emph{Proceedings of the National Academy of Sciences} \textbf{111}(9), 3316–3321 (2014).
		
		\bibitem{Jenkyns2010}
		H. C. Jenkyns, ``Geochemistry of oceanic anoxic events,''
		\emph{Geochemistry, Geophysics, Geosystems} \textbf{11}(3), Q03004 (2010).
		
		\bibitem{OAE2Timing}
		S. Kolonic, T. Wagner, A. Forster, J. S. Sinninghe Damsté, ``Black shale deposition on the northwest African Shelf during the Cenomanian/Turonian oceanic anoxic event: Climate coupling and global organic carbon burial,''
		\emph{Paleoceanography} \textbf{20}(1), PA1006 (2005).
		
		\bibitem{Laschamp}
		Q. Simon, N. Thouveny, D. L. Bourlès, J. Valet, F. Bassinot, ``Cosmogenic \(^{10}\)Be production records reveal dynamics of geomagnetic dipole moment (GDM) over the Laschamp excursion (20–60 ka),''
		\emph{Earth and Planetary Science Letters} \textbf{550}, 116547 (2020).
		
		\bibitem{Jablonski1991}
		D. Jablonski, ``Extinctions: A paleontological perspective,''
		\emph{Science} \textbf{253}(5021), 754–757 (1991).
		
		\bibitem{CampbellAllen2008}
		I. H. Campbell, C. M. Allen, ``Supermountains and the rise of atmospheric oxygen,''
		\emph{Nature Geoscience} \textbf{1}, 554 (2008).
		
		\bibitem{RampinoCaldeira2015}
		M. R. Rampino, K. Caldeira, ``Periodic impact cratering and extinction events over the last 260 million years,''
		\emph{Monthly Notices of the Royal Astronomical Society} \textbf{454}(4), 3480–3484 (2015). DOI: 10.1093/mnras/stv2088
		
		\bibitem{MateseWhitmire2011}
		J. J. Matese, D. P. Whitmire, ``Persistent evidence of a jovian mass solar companion in the Oort cloud,''
		\emph{Icarus} \textbf{211}(2), 926–938 (2011).
		
		\bibitem{Whitmire2025}
		D. P. Whitmire, J. J. Matese, ``New constraints on the Nemesis hypothesis from the WISE survey,''
		\emph{Icarus} \textbf{400}, 115654 (2025).
		
		\bibitem{Crumpler2024}
		N. R. Crumpler, V. Chandra, N. L. Zakamska, ``Detection of the temperature dependence of the white dwarf mass-radius relation with gravitational redshifts,''
		\emph{The Astrophysical Journal} \textbf{977}, 237 (2024). DOI: 10.3847/1538-4357/ad8ddc
		
		\bibitem{Lantink2022}
		M. L. Lantink, J. H. F. L. Davies, M. Ovtcharova, F. J. Hilgen, ``Milankovitch cycles in banded iron formations constrain the Earth-Moon system 2.46 billion years ago,''
		\emph{Nature Geoscience} \textbf{15}, 571–576 (2022).
		
		\bibitem{Farhat2022}
		M. Farhat, P. Auclair-Desrotour, G. Boué, J. Laskar, ``The resonant tidal evolution of the Earth-Moon distance,''
		\emph{Astronomy \& Astrophysics} \textbf{665}, A108 (2022).
		
		\bibitem{Rosat2025}
		S. Rosat, C. Gaugne Gouranton, I. Panet, M. Mandea, ``GRACE detection of transient mass redistributions during a mineral phase transition in the deep mantle,''
		\emph{Geophysical Research Letters} \textbf{52}, e2025GL116408 (2025). DOI: 10.1029/2025GL116408
		
		\bibitem{Planck2020}
		Planck Collaboration, ``Planck 2018 results. VI. Cosmological parameters,''
		\emph{Astronomy \& Astrophysics} \textbf{641}, A6 (2020).
		
		\bibitem{Nesvorny2002}
		D. Nesvorný, W. F. Bottke, L. Dones, H. F. Levison, ``The recent breakup of an asteroid in the main-belt region,''
		\emph{Nature} \textbf{417}, 720–722 (2002).
		
		\bibitem{Tamayo2020}
		D. Tamayo, H. Rein, P. Shi, D. M. Hernandez, ``ASSIST: An ephemeris-quality test-particle integrator,''
		\emph{The Planetary Science Journal} \textbf{4}, 69 (2020).
		
		\bibitem{Rehren2013}
		T. Rehren, L. Belmonte, A. Pankhurst, L. Schiavon, D. A. Scott, J. W. Taylor, ``5,000 years old Egyptian iron beads from Gerzeh were made from hammered meteoritic iron,''
		\emph{Journal of Archaeological Science} \textbf{40}, 4785–4792 (2013).
		
		\bibitem{Hofmann2023}
		B. Hofmann, S. Schreyer, S. A. Sorensen, F. R. R. P. de Meijer, ``The Mörigen arrowhead — an iron meteorite found in a Bronze Age lake dwelling,''
		\emph{Journal of Archaeological Science} \textbf{156}, 105800 (2023).
		
		\bibitem{JutziAsphaug2011}
		M. Jutzi, E. Asphaug, ``The Moon's far side dichotomy explained by a low-velocity collision with a companion moon,''
		\emph{Nature} \textbf{476}, 69–72 (2011).
		
		\bibitem{Gardiner2025}
		A. A. Gardiner, R. J. Walker, A. L. D. Kilburn, ``Native hydrogen in enstatite chondrite LAR 12252: Implications for Earth's water delivery,''
		\emph{Icarus} (accepted, 2025).
		
		\bibitem{Rampino2017}
		M. R. Rampino, \emph{Cataclysms: A New Geology for the Twenty-First Century},
		Columbia University Press (2017).
		
		\bibitem{Yakushev2026}
		A. V. Yakushev, ``Yakushev Unified Coordination Theory (YUCT),''
		Zenodo, \url{https://doi.org/10.5281/zenodo.18444599} (2026).
		
		\bibitem{AppendixL}
		A. V. Yakushev, ``Appendix L: Fractal Coordination Error Scaling — A Universal Law from DNA to Cosmology,''
		in \emph{YUCT}, Zenodo (2026).
		
		\bibitem{AppendixP}
		A. V. Yakushev, ``Appendix P: Temperature Dependence of Gravity — Observational Confirmation and Theoretical Foundation within the Coordination Paradigm (YUCT),''
		in \emph{YUCT}, Zenodo (2026).
		
		\bibitem{AppendixW}
		A. V. Yakushev, ``Appendix W: Passive Gravitational Tomography of Planetary Interiors Using Tidal Waves — A Coordination-Based Approach,''
		in \emph{YUCT}, Zenodo (2026).
		
		\bibitem{AppendixY}
		A. V. Yakushev, ``Appendix Y: Mathematical Foundations of YUCT — A Derivation from First Principles,''
		in \emph{YUCT}, Zenodo (2026).
		
		\bibitem{AppendixΩ}
		A. V. Yakushev, ``Appendix Ω: The Global Rotation of the Universe and Its New Minimum Age from the Dipole Turning Radius,''
		in \emph{YUCT}, Zenodo (2026).
		
		\bibitem{Whitmire1985}
		D. P. Whitmire, J. J. Matese, ``Periodic comet showers and the Nemesis hypothesis,''
		\emph{Nature} \textbf{313}, 36–38 (1985).
		
		\bibitem{Kirkpatrick2021}
		D. Kirkpatrick et al., ``The field brown dwarf population and the WISE survey,''
		\emph{Astrophysical Journal Supplement} \textbf{253}, 7 (2021).
		
		\bibitem{Korenaga2017}
		J. Korenaga, ``Earth’s thermal evolution, mantle convection, and Hadean onset of plate tectonics,''
		\emph{Earth and Planetary Science Letters} \textbf{145}, 334–348 (2017).
		
		\bibitem{Charnay2017}
		B. Charnay, G. Le Hir, F. Fluteau, F. Forget, D. C. Catling, ``A warm or a cold early Earth? New insights from a 3‑D climate‑carbon model,''
		\emph{Earth and Planetary Science Letters} \textbf{474}, 97–109 (2017).
		
		\bibitem{KrissansenTotton2019}
		J. Krissansen‑Totton, G. Arney, D. C. Catling, ``Ocean pH, atmospheric CO₂, and habitability of the early Earth,''
		\emph{AGU Fall Meeting} PP53B‑06 (2019).
		
		\bibitem{Arndt2012}
		N. T. Arndt, E. G. Nisbet, ``Processes on the Young Earth and the Habitats of Early Life,''
		\emph{Annual Review of Earth and Planetary Sciences} \textbf{40}, 521–549 (2012).
		
		\bibitem{Zircon2024}
		H. Gamaleldien et al., ``Four‑billion‑year‑old zircons reveal early fresh water and emerged continental crust,''
		\emph{Nature Geoscience} (2024). DOI: 10.1038/s41561‑024‑01450‑0
		
		\bibitem{PanspermiaReview}
		P. A. Bland, ``Panspermia: A viable hypothesis?''
		\emph{Astronomy \& Geophysics} \textbf{58}(5), 5.24–5.28 (2017).
		
		\bibitem{LunarMagmaOcean}
		M. Maurice, N. Tosi, S. Schwinger, D. Breuer, T. Kleine, ``The longevity of a lunar magma ocean: Implications for the Moon's thermal and magnetic evolution,''
		\emph{Journal of Geophysical Research: Planets} \textbf{125}, e2019JE006152 (2020). DOI: 10.1029/2019JE006152
		
		\bibitem{Chemolithoautotroph}
		Wikipedia, ``Chemolithoautotroph,''
		\url{https://en.wikipedia.org/wiki/Chemolithoautotroph} (accessed 2026).
		
		\bibitem{AbiogenesisReview}
		N. Kitadai, S. Maruyama, ``Onset of plate tectonics and the origin of life: A view from the Hadean,''
		\emph{Geoscience Frontiers} \textbf{9}(4), 1105–1122 (2018). DOI: 10.1016/j.gsf.2017.05.005
		
		\bibitem{Chemoautotroph}
		K. O. Stetter, ``Hyperthermophilic procaryotes,''
		\emph{FEMS Microbiology Reviews} \textbf{18}(2-3), 149–158 (1996).
		
		\bibitem{JanusEpimetheus}
		F. Namouni, H. Christophe, ``On the stability of co‑orbital motion of Janus and Epimetheus,''
		\emph{Astronomy \& Astrophysics} \textbf{434}, 1179–1186 (2005).
		
		\bibitem{Kearey2013}
		P. Kearey, K. A. Klepeis, F. J. Vine, \emph{Global Tectonics}, 3rd ed., Wiley-Blackwell (2013).
		
		\bibitem{Sahijpal2020}
		S. Sahijpal, V. Goyal, ``Thermal evolution of the early Moon,''
		\emph{Meteoritics \& Planetary Science} \textbf{53}(10), 2193–2210 (2018). DOI: 10.1111/maps.13119
		
		\bibitem{Colin2024}
		L. Colin, M. Maurice, N. Tosi, et al., ``Thermal evolution of the lunar magma ocean,''
		\emph{Earth and Planetary Science Letters} \textbf{648}, 119072 (2024). DOI: 10.1016/j.epsl.2024.119072
		
		\bibitem{Maurice2018}
		M. Maurice, N. Tosi, S. Schwinger, et al., ``A long-lived lunar magma ocean,''
		\emph{AGU Fall Meeting Abstracts}, P43C-06 (2018).
		
		\bibitem{Maurice2020}
		M. Maurice, N. Tosi, S. Schwinger, et al., ``A long-lived magma ocean on a young Moon,''
		\emph{Science Advances} \textbf{6}(28), eaba8949 (2020). DOI: 10.1126/sciadv.aba8949
		
		\bibitem{Zahnle2015}
		K. J. Zahnle, R. Lupu, D. C. Catling, ``The evolution of a steam atmosphere after a giant impact,''
		\emph{AGU Fall Meeting} P51A-1920 (2015).
		
		\bibitem{Zahnle2007}
		K. J. Zahnle, N. Arndt, C. Cockell, et al., ``Emergence of a habitable planet,''
		\emph{Space Science Reviews} \textbf{129}(1-3), 35–78 (2007). DOI: 10.1007/s11214-007-9225-z
		
		\bibitem{Feulner2012}
		G. Feulner, ``The faint young Sun problem,''
		\emph{Reviews of Geophysics} \textbf{50}(2), RG2006 (2012). DOI: 10.1029/2011RG000375
		
		\bibitem{Canup2014}
		R. M. Canup, ``Lunar-forming impacts: processes and alternatives,''
		\emph{Philosophical Transactions of the Royal Society A} \textbf{372}(2024), 20130175 (2014). DOI: 10.1098/rsta.2013.0175
		
	\end{thebibliography}
	
\end{document}